\documentclass[12pt,a4paper]{report}

\usepackage[spanish]{babel}
\usepackage[utf8]{inputenc}
\usepackage[T1]{fontenc}
\usepackage{lmodern}
\usepackage{amsmath,amssymb}
\usepackage{graphicx}
\usepackage{booktabs}
\usepackage{enumitem}
\usepackage{geometry}
\usepackage{titlesec}
\usepackage{csquotes}
\usepackage{hyperref}
\usepackage{setspace}
\usepackage[
  backend=biber,
  style=authoryear,
  bibstyle=numeric,
  backref=true,
  maxcitenames=1,
  maxbibnames=99,
  hyperref=true
]{biblatex}
% \usepackage{biblatex}

\geometry{margin=3.5cm}
\emergencystretch=2em
\addbibresource{memoria.bib}
\setstretch{1.5}

\titlespacing*{\chapter}{0pt}{80pt}{50pt}
\titlespacing*{\section}{0pt}{36pt}{18pt}
\titlespacing*{\subsection}{0pt}{24pt}{12pt}

\newcommand{\citationneeded}[1]{\textbf{[Falta cita: #1]}}
\newcommand{\verifycite}[1]{\textbf{[Verificar: #1]}}

%%%%%%%%%%%%%%%%%%%%%%%%%%

\title{Aprendizaje por refuerzo para la ciberseguridad: un defensor basado en flujos para decisiones binarias de permitir/bloquear}
\author{Javier Rivero Iglesias}
\date{Mayo de 2026}

%%%%%%%%%%%%%%%%%%%%%%%%%%

\begin{document}

\pagenumbering{gobble}
\maketitle
\newpage
\tableofcontents
\newpage
\pagenumbering{arabic}

\chapter{Introducción}

Este capítulo sitúa el problema de investigación, delimita una especificación verificable para el prototipo y adelanta las contribuciones, el alcance y la organización de la memoria. La intención es distinguir desde el inicio entre la evidencia que puede obtenerse mediante experimentación fuera de línea (\emph{offline}) y cualquier afirmación de preparación para un despliegue operativo.

\section{Motivación}
Las infraestructuras de red modernas generan volúmenes de tráfico que superan con creces lo que los analistas humanos pueden inspeccionar manualmente. Una única red empresarial puede producir millones de registros de flujo al día, y los despliegues a gran escala agregan flujos procedentes de miles de endpoints, servicios en la nube y enlaces entre centros de datos \parencite{ScarfoneMell2007}. Distinguir la comunicación legítima de la actividad maliciosa dentro de este volumen es uno de los desafíos operativos centrales en la seguridad de redes. Las respuestas tradicionales se apoyan en sistemas de detección de intrusiones en red (NIDS, del inglés \emph{Network Intrusion Detection Systems}) que aplican coincidencia de firmas, umbrales estadísticos o reglas específicas de protocolo para señalar actividad sospechosa. Estos enfoques son efectivos frente a patrones de amenaza conocidos, pero requieren actualizaciones manuales constantes de reglas, resultan en gran medida ineficaces frente a variantes novedosas de ataque y generan volúmenes de alertas que habitualmente desbordan a los equipos de seguridad.

El aprendizaje automático (ML, del inglés \emph{Machine Learning}) ofrece una perspectiva complementaria. En lugar de codificar el conocimiento experto como reglas estáticas, los modelos de ML pueden aprender patrones estadísticos a partir de datos de tráfico etiquetados y generalizar hacia instancias de ataque no vistas previamente. A medida que el campo ha madurado, las representaciones a nivel de flujo se han convertido en un sustrato práctico para este aprendizaje: cada conexión de red se resume como un conjunto fijo de características estadísticas que cubren duración, conteos de paquetes, totales de bytes, tiempos entre llegadas, distribuciones de flags TCP y métricas basadas en tasas \parencite{Sperotto2010FlowIDS,Umer2017FlowIDS}. Esta reducción hace abordable el aprendizaje a gran escala y evita las complicaciones legales y técnicas de la inspección profunda de paquetes, que resulta cada vez más ineficaz de todos modos a medida que crece la fracción de tráfico cifrado en internet.

La limitación clave del aprendizaje supervisado estándar en este contexto es que trata la detección de intrusiones como un problema de clasificación simétrico, minimizando el error de clasificación global sin tener en cuenta las consecuencias operativas de los distintos tipos de error \parencite{Axelsson1999BaseRate}. En la práctica, esas consecuencias son profundamente asimétricas. Un falso negativo permite a un actor de amenazas avanzar sin ser detectado y puede conducir a exfiltración de datos, despliegue de ransomware o interrupción del servicio. Un falso positivo genera una alerta o un bloqueo de tráfico que consume tiempo del analista, pero rara vez causa daños duraderos. El aprendizaje por refuerzo (RL, del inglés \emph{Reinforcement Learning}) proporciona un marco natural para codificar esta asimetría directamente en el objetivo de aprendizaje: al diseñar una función de recompensa que penaliza los ataques no detectados más severamente que las falsas alarmas, un agente de RL puede ser guiado hacia políticas que reflejen las prioridades de riesgo operativo real en lugar de la exactitud agregada \parencite{SuttonBarto2018RL,Lee2002CostSensitiveIDS}. Explorar cómo un agente basado en la \emph{Quantile Regression Deep Q-Network} (QRDQN) navega estos compromisos sobre datos a nivel de flujo es la motivación central de esta tesis \parencite{Dabney2018QRDQN}.

\section{Definición del problema}

\subsection{Desafíos metodológicos}
Aplicar el aprendizaje por refuerzo a la detección de intrusiones en red plantea varios desafíos metodológicos que van más allá de sustituir un algoritmo de RL por un clasificador estándar. El primer desafío es la calidad de los datos. Los bancos de pruebas públicos (\emph{benchmarks}) de NIDS como CICIDS2017 son ampliamente utilizados, pero contienen artefactos conocidos de su flujo de procesamiento (\emph{pipeline}) de extracción de flujos, entre ellos ruido en las etiquetas, características duplicadas, fallos de implementación de CICFlowMeter y campos estadísticos que codifican inadvertidamente el contexto del día de captura \parencite{Ring2019DatasetSurvey,Engelen2021CICIDSIssues,Lanvin2023Errors}. Los modelos entrenados sin auditar estos artefactos pueden aprender a reconocer firmas de extracción o identificadores específicos de escenario en lugar de comportamiento de ataque genuino, produciendo puntuaciones infladas en el \emph{benchmark} que no se transfieren a ningún entorno nuevo.

El segundo desafío es la fuga de información. Muchos estudios publicados de NIDS evalúan los modelos utilizando particiones aleatorias por filas sobre todo el tráfico disponible, lo que puede situar flujos altamente correlacionados a ambos lados del límite entrenamiento-prueba. Más allá de las filas duplicadas, la fuga puede introducirse a través de direcciones IP, marcas temporales absolutas, identificadores del flujo \textit{(Flow IDs)} y números de puerto que son específicos de un escenario de ataque programado en lugar de indicativos del comportamiento malicioso en general \parencite{Arp2020DosDontsMLSecurity,Kaufman2011LeakageDataMining}. Los pasos de preprocesamiento como el escalado de características, aplicados globalmente antes de la partición, pueden contaminar adicionalmente la evaluación al permitir que la información distribucional de la partición de prueba influya en la normalización del entrenamiento. Un estudio riguroso no puede limitarse a elegir un algoritmo de alto rendimiento y reportar su exactitud sobre un conjunto de datos público dividido aleatoriamente.

El tercer desafío es la comparación justa contra una línea base (\emph{baseline}) de un método estándar. Los métodos de aprendizaje por refuerzo profundo (DRL, del inglés \emph{Deep Reinforcement Learning}) pueden ser costosos de ajustar, sensibles a los hiperparámetros y más lentos en converger que los modelos supervisados tabulares. El bosque aleatorio (\emph{Random Forest}) y los ensamblados de árboles con potenciación de gradiente (\emph{gradient boosting}) son líneas base sólidas para la detección de intrusiones a nivel de flujo y pueden alcanzar puntuaciones casi perfectas en CICIDS2017 bajo condiciones de evaluación favorables \parencite{LiuLang2019IDSSurvey,Apruzzese2022CrossEvaluationNIDS}. Un defensor basado en RL no puede evaluarse de manera significativa sin una comparación bajo el mismo protocolo frente a dichas líneas base. La pregunta no es si QRDQN alcanza alta exactitud en una sola ejecución, sino si su formulación sensible al coste ofrece alguna ventaja empírica sobre un modelo supervisado bien ajustado bajo esquemas de características, protocolos de partición y métricas de evaluación equivalentes.

\subsection{Especificación verificable del proyecto}
\label{sec:especificacion-proyecto}
Como adaptación al contexto de investigación de la especificación propia de un proyecto de desarrollo de software (PDS), esta tesis formula el trabajo como un sistema experimental verificable y no como un producto de bloqueo listo para operar. La Tabla~\ref{tab:especificacion-proyecto} fija el problema que se aborda, los límites que preservan la seguridad metodológica y la evidencia exigida para considerar cubierto cada compromiso.

\begin{table}[htbp]
\centering
\small
\begin{tabular}{@{}p{2.6cm}p{9.7cm}@{}}
\toprule
\multicolumn{2}{@{}p{12.3cm}@{}}{\textbf{Especificación verificable del prototipo experimental}} \\
\midrule
Necesidad & Construir y evaluar un defensor offline que clasifique flujos de red como \texttt{PERMIT} o \texttt{BLOCK}, haciendo explícita la asimetría entre permitir un ataque y bloquear tráfico legítimo. \\
Entradas & CSV de flujos de CICIDS2017 y, para la comprobación complementaria, flujos extraídos de tráfico de laboratorio; todos se adaptan al contrato canónico y se procesan sin usar identificadores ni otros campos propensos a fuga. \\
Proceso & Limpieza y escalado ajustados solo con entrenamiento, representación de $76+76=152$ dimensiones, entrenamiento QRDQN, comparación con Random Forest y escalera de validación con artefactos persistidos. \\
Salidas & Predicciones binarias experimentales, métricas por clase y artefactos de ejecución trazables; \texttt{BLOCK} permanece como salida de inferencia offline y no acciona un cortafuegos ni una red real. \\
Verificación & Los compromisos C1--C4, especificados en las Tablas~\ref{tab:obj-c1} a~\ref{tab:obj-c4}, requieren código auditable, ejecuciones comprometidas y resultados interpretados dentro de las restricciones declaradas. \\
Límite de la afirmación & El rendimiento sobre CICIDS2017 y sobre una captura concreta de laboratorio constituye evidencia experimental acotada; no certifica generalización a producción ni bloqueo activo seguro. \\
\bottomrule
\end{tabular}
\caption[Especificación verificable del proyecto]{Especificación verificable del prototipo experimental, adaptada al contexto de un proyecto de desarrollo de software. Fuente: elaboración propia.}
\label{tab:especificacion-proyecto}
\end{table}

\section{Enfoque propuesto}
Esta tesis formula la tarea defensiva como un problema de aprendizaje por refuerzo offline con el conjunto de datos como entorno \parencite{Levine2020OfflineRL}. La elección del entorno offline es deliberada: desplegar un agente de RL que deba interactuar con una red en vivo para aprender generaría riesgos de seguridad inaceptables, ya que un agente en exploración tomaría inevitablemente acciones subóptimas durante el entrenamiento. Al tratar un conjunto de datos etiquetado estático como el entorno, el agente puede entrenarse y evaluarse sin ninguna exposición a una red en vivo. El tráfico de red se procesa en un esquema canónico de flujo de 76 características fijo, independiente del formato nativo de cualquier conjunto de datos específico. Este contrato canónico separa la nomenclatura de columnas propia de la extracción del espacio de observación del modelo y hace posible procesar tanto flujos de CICIDS2017 como tráfico capturado en un laboratorio doméstico cerrado a través del mismo pipeline sin modificación alguna.

Para gestionar esquemas heterogéneos al transitar entre herramientas de extracción, se añade una máscara de presencia/ausencia de 76 valores, produciendo una observación de 152 dimensiones. Cada posición indica si existe una columna fuente mapeada para la característica canónica correspondiente. En CICIDS2017 las 76 columnas canónicas están cubiertas y la máscara es constante; en la Fase~2 puede revelar columnas no disponibles. Los valores no finitos de una columna existente se sanean antes del mapeo y no cambian su indicador. Esta semántica, más estrecha que una máscara de validez fila a fila, coincide con la implementación actual y hace auditable la compatibilidad estructural entre extractores.

La tarea de decisión se mapea a un espacio de acciones binario: \texttt{PERMIT} para el tráfico clasificado como benigno, y \texttt{BLOCK} para el tráfico clasificado como ataque. Un agente \textit{Quantile Regression Deep Q-Network} \parencite{Dabney2018QRDQN}, implementado a través de la biblioteca \texttt{sb3-contrib} \parencite{SB3ContribQRDQNDocs}, se entrena dentro de un entorno compatible con Gymnasium \parencite{GymnasiumDocs}. El entorno aplica una función de recompensa asimétrica que penaliza los falsos negativos más severamente que los falsos positivos, codificando directamente las prioridades de riesgo operativo descritas anteriormente \parencite{Cardenas2006IDSFramework}. El conjunto de datos CICIDS2017 \parencite{Sharafaldin2018CICIDS2017,CICIDS2017Web} sirve como benchmark interno primario, procesado a través de un pipeline de carga curado que excluye los campos propensos a fuga de información antes de que comience cualquier entrenamiento del modelo. El rendimiento se evalúa frente a una línea base de Random Forest bajo el mismo esquema de características y protocolo de partición \parencite{Engelen2021CICIDSIssues}. Un segundo nivel de validación extiende la evaluación ejecutando inferencia offline sobre características de flujo extraídas de tráfico capturado por el operador en un laboratorio doméstico cerrado y no adversarial, ofreciendo una indicación parcial y de validez externa limitada del comportamiento del sistema ante un cambio de dominio.

\section{Contribución de la tesis}
Las contribuciones se expresan mediante los cuatro compromisos contractuales del Capítulo~2. Esta formulación permite separar el producto construido, la evidencia de evaluación y los límites de cada conclusión, en lugar de presentar el rendimiento de una sola partición como una demostración global.

\textbf{C1 --- Pipeline experimental reproducible con contrato de datos canónico.} La primera contribución es un pipeline offline de extremo a extremo que fija una representación de $76+76=152$ dimensiones y conserva los artefactos necesarios para reconstruir cada resultado. La Tabla~\ref{tab:obj-c1} especifica este compromiso y el Anexo~\ref{cap:anexo-artefactos} documenta el contrato de evidencia.

\textbf{C2 --- Defensor RL y jerarquía de validación con control de fugas.} La segunda contribución es la implementación de QRDQN como defensor binario y el diseño de una campaña que separa MAIN aleatoria, día completo, tamaño, semilla y cuatro escenarios retenidos, junto con controles directos, barajado de etiquetas y duplicados. La Tabla~\ref{tab:obj-c2} fija el compromiso; el Capítulo~\ref{cap:protocolo-experimental} distingue el diseño aprobado de la evidencia todavía pendiente.

\textbf{C3 --- Decisión sensible al coste con lectura multiobjetivo.} La tercera contribución convierte la asimetría entre ataques permitidos y tráfico legítimo bloqueado en un elemento auditable del diseño, no en una interpretación posterior de la exactitud. La Tabla~\ref{tab:obj-c3} enlaza esta contribución con la matriz de recompensa y su lectura de seguridad e impacto operativo (Tabla~\ref{tab:recompensa-multiobjetivo}).

\textbf{C4 --- Comparación honesta y transferencia offline acotada.} La cuarta contribución incorpora Random Forest bajo particiones emparejadas y una Fase~2 que reutiliza modelo y preprocesamiento sin reajuste. La evidencia histórica se conserva como antecedente; la comparación final y la inferencia fresca dependen de la campaña. La Tabla~\ref{tab:obj-c4} fija el compromiso y la Tabla~\ref{tab:estructura-resultados-finales} define la evidencia necesaria, sin atribuir a \texttt{BLOCK} ningún efecto operativo.

\section{Alcance y limitaciones}
El alcance de esta investigación se limita estrictamente a la evaluación offline, tabular y a nivel de flujo. Las salidas \texttt{PERMIT} y \texttt{BLOCK} son clasificaciones binarias experimentales sobre registros de flujo estáticos; el sistema no realiza bloqueo en vivo de cortafuegos, descarte de paquetes, manipulación de \texttt{iptables} ni prevención de intrusiones en línea. El entorno no simula un adversario activo cuyas acciones posteriores dependan de las decisiones del defensor, lo que significa que el sistema no opera como un plano de control de toma de decisiones secuenciales completo. Estas restricciones son decisiones de diseño deliberadas que hacen que la evaluación sea abordable, reproducible y segura. No son omisiones que deban subsanarse en iteraciones futuras de este mismo prototipo.

CICIDS2017 se utiliza como benchmark interno primario para la experimentación controlada. Los resultados sobre este conjunto de datos deben interpretarse como evidencia sobre el comportamiento bajo un benchmark específico y bien caracterizado, no como prueba de generalización a redes de producción arbitrarias. El tráfico de laboratorio, generado por el operador en un entorno doméstico cerrado y no adversarial, se utiliza como nivel de evaluación secundario para examinar el cambio de dominio; dada su procedencia, su validez externa es limitada y, además, también es inferencia offline y no despliegue en vivo. La tesis ocupa una posición metodológica entre la evaluación pura sobre benchmark y la validación de despliegue completo, apuntando a la reproducibilidad e interpretabilidad del marco experimental más que a la operatividad en producción.

\section{Estructura del documento}
Tras esta introducción, el Capítulo~2 fija los objetivos contractuales y específicos, el alcance, las restricciones y la planificación. El Capítulo~3 revisa el estado del arte de los NIDS basados en flujos, los conjuntos de datos públicos, el aprendizaje automático y el aprendizaje por refuerzo aplicados a detección de intrusiones, junto con sus riesgos metodológicos. El Capítulo~4 presenta el diseño del sistema: datos, limpieza anti-fuga, contrato canónico, formulación del entorno y agente QRDQN. El Capítulo~5 define el protocolo experimental, las particiones, las métricas, la línea base y la inferencia offline de la Fase~2.

El Capítulo~6 fija la estructura de incorporación de la campaña final y conserva, claramente separada, la evidencia histórica previa. El Capítulo~7 organiza la discusión por generalización, duplicados, coste, comparación de modelos, escala, semillas, escenarios y cambio de dominio; el Capítulo~8 expone las limitaciones y el Capítulo~9 delimita las implicaciones éticas. El Capítulo~10 presenta una conclusión provisional que podrá cerrarse desde los artefactos finales. Los anexos reúnen reproducción, entorno, esquema canónico y catálogo de artefactos.

\chapter{Objetivos y alcance}

\section{Objetivo general}
El objetivo general de esta tesis es diseñar, implementar y evaluar rigurosamente un defensor de ciberseguridad basado en aprendizaje por refuerzo, de carácter offline, que toma decisiones binarias sensibles al coste (\texttt{PERMIT} o \texttt{BLOCK}) sobre observaciones de flujos de red estandarizados. Este objetivo se persigue mediante un \textit{pipeline} experimental controlado que opera íntegramente sobre datos de flujo etiquetados y estáticos, priorizando el rigor metodológico y la reproducibilidad sobre cualquier pretensión de preparación para el despliegue.

El enfoque se fundamenta en la observación de que la detección de intrusiones implica costes inherentemente asimétricos: permitir un ataque conlleva consecuencias que son, en general, considerablemente más graves que las de generar una falsa alarma. Una formulación basada en recompensas hace explícita esta asimetría en el objetivo de aprendizaje, en lugar de dejarla implícita en un ajuste posterior de umbrales. La pregunta central que guía el trabajo es si un algoritmo de RL distribucional como QRDQN puede configurarse, entrenarse y evaluarse de forma que demuestre un comportamiento sensible al coste y medible sobre un benchmark público bien estudiado, manteniéndose a la vez comparable con sólidas líneas base supervisadas bajo las mismas condiciones controladas.

\section{Objetivos específicos}

\subsection{Representación del tráfico basada en flujos}
El primer objetivo específico es extraer y normalizar datos de tráfico de red en registros tabulares de flujos, garantizando que el formato de entrada sea compatible tanto con benchmarks públicos como CICIDS2017 como con capturas de laboratorio offline. Los flujos de red son resúmenes bidireccionales de intercambios de paquetes que comparten una quíntupla común, y su estructura tabular compacta y de anchura fija los hace muy adecuados como entradas para modelos de aprendizaje automático \parencite{Sperotto2010FlowIDS,Umer2017FlowIDS}. Alcanzar este objetivo requiere auditar las características disponibles en los ficheros CSV generados por CICFlowMeter, identificar qué campos contienen señal estadística legítima y cuáles representan artefactos de extracción o posibles vectores de fuga de información.

Este objetivo implica también definir un contrato de características estable que tienda un puente entre las publicaciones académicas en formato CSV y las exportaciones brutas de CICFlowMeter. Un contrato estable significa que el conjunto de características utilizado como entrada al modelo no varía entre experimentos, que los campos excluidos por razones de fuga se excluyen de forma consistente, y que cualquier persona que lea el código fuente puede verificar exactamente qué medidas llegan al modelo. Esta estabilidad es un requisito previo para todos los objetivos posteriores, ya que todos los demás componentes del pipeline dependen de un espacio de observación bien definido y auditable.

\subsection{Preprocesamiento canónico y esquema de características}
El segundo objetivo específico es definir y hacer cumplir un esquema canónico estricto de 76 características junto con una máscara de presencia/ausencia de 76 valores, lo que da lugar a un vector de observación estable de 152 dimensiones. Las 76 características canónicas se extraen de estadísticas de estilo CICFlowMeter que cubren la duración del flujo, conteos de paquetes y bytes, distribuciones de tiempos de llegada, conteos de flags TCP, estadísticas de longitud de paquetes, resúmenes de periodos activos e inactivos, y métricas de tasa de flujo. Los campos que podrían exponer información de identidad específica del conjunto de datos, como las direcciones IP de origen y destino, las marcas temporales absolutas, los \textit{Flow ID} y los números de puerto que funcionan como proxies directos del tipo de ataque, se excluyen explícitamente a nivel de esquema, en lugar de filtrarse en pasos de preprocesamiento ad hoc.

La máscara de presencia/ausencia de 76 valores registra si cada valor canónico fue medido directamente y está presente (codificado como 1) o ausente y sustituido por un marcador de posición por defecto (codificado como 0). Este diseño hace visibles las lagunas de datos tanto para el modelo como para el evaluador, evitando una situación en la que la red se entrena silenciosamente sobre ceros imputados sin ninguna indicación de que la medición subyacente no estaba disponible. Al pasar de CICIDS2017 a tráfico capturado en laboratorio, distintos pipelines de extracción pueden producir conjuntos de columnas diferentes; la máscara de presencia/ausencia garantiza que cualquier discrepancia de este tipo quede codificada en la observación en lugar de quedar oculta, de modo que el lector pueda distinguir un cero medido de uno imputado. Juntos, el esquema canónico y la máscara constituyen el contrato de observación del que depende cada componente posterior.

\subsection{Entorno de aprendizaje por refuerzo}
El tercer objetivo específico es construir una interfaz conjunto-de-datos-como-entorno compatible con Gymnasium \parencite{GymnasiumDocs} que presente los registros de flujo al agente como una secuencia de observaciones. El entorno envuelve un conjunto de datos etiquetado estático de modo que cada llamada a \texttt{step()} devuelve un registro de flujo como observación, junto con la recompensa calculada a partir de la acción anterior del agente y la etiqueta verdadera del último registro. Esta formulación permite que los bucles de entrenamiento de RL estándar, desarrollados para entornos interactivos, se apliquen directamente a conjuntos de datos de flujo offline sin modificar el algoritmo.

El entorno debe implementar una función de recompensa asimétrica y sensible al coste que penalice los falsos negativos más que los falsos positivos, reflejando los perfiles de riesgo en ciberseguridad en los que los ataques no detectados son operacionalmente más perjudiciales que las falsas alarmas \parencite{Lee2002CostSensitiveIDS,Cardenas2006IDSFramework}. Los pesos de recompensa específicos se tratan como parámetros experimentales configurables en lugar de constantes universales fijas, ya que la relación de costes adecuada depende del modelo de amenaza y del contexto operativo del escenario de despliegue. El entorno debe también admitir reinicios reproducibles, gestión de límites de episodios y particiones del conjunto de datos configurables, de modo que la misma interfaz pueda utilizarse tanto para el entrenamiento como para la evaluación bajo distintos modos de partición.

\subsection{Agente de decisión binaria basado en QRDQN}
El cuarto objetivo específico es configurar y entrenar un agente Quantile Regression Deep Q-Network \parencite{Dabney2018QRDQN}, aprovechando su enfoque distribucional para la estimación del valor, para navegar el espacio de acciones binario \texttt{PERMIT}/\texttt{BLOCK} a partir de las observaciones canónicas de 152 dimensiones. QRDQN extiende el DQN estándar estimando la distribución completa de los retornos esperados en lugar de solo la media, utilizando un conjunto de cuantiles aprendidos y una pérdida de regresión cuantílica. Esta representación distribucional resulta de interés en un contexto sensible al coste porque proporciona información más rica sobre la dispersión y la forma de los posibles resultados, no solo el retorno promedio, y en principio permite que la política sea más conservadora ante los peores casos.

La implementación práctica utiliza el módulo QRDQN de la biblioteca \texttt{sb3-contrib} \parencite{SB3ContribQRDQNDocs}, que proporciona una implementación estable y probada del algoritmo con arquitectura de red configurable, buffer de repetición, esquema de exploración e hiperparámetros de entrenamiento. El objetivo incluye seleccionar una configuración razonable por defecto, validar que el bucle de entrenamiento converge en el benchmark interno, y persistir todos los artefactos del modelo entrenado junto con la configuración y el estado de preprocesamiento utilizados para producirlos. La reproducibilidad exige que un modelo almacenado pueda cargarse y evaluarse de forma idéntica en una fecha posterior sin necesidad de reconstruir el pipeline de entrenamiento completo desde cero.

\subsection{Comparación con línea base supervisada}
El quinto objetivo específico es implementar una línea base supervisada robusta, utilizando principalmente Random Forest \parencite{LiuLang2019IDSSurvey}, para ofrecer una evaluación comparativa justa, con el mismo protocolo, del valor empírico del agente QRDQN. Random Forest se elige como línea base principal porque está bien establecido para la clasificación tabular, gestiona interacciones no lineales entre características sin un preprocesamiento extenso, es robusto ante características irrelevantes y tiende a obtener buenos resultados en benchmarks de NIDS basados en flujos \parencite{Apruzzese2022CrossEvaluationNIDS}. Utilizar el mismo esquema de características, el mismo protocolo de partición y las mismas métricas de evaluación tanto para el agente de RL como para la línea base es esencial; cualquier discrepancia en el preprocesamiento, la disponibilidad de características o las condiciones de evaluación comprometería la validez de la comparación.

Esta comparación cumple varios propósitos. En primer lugar, determina si la formulación QRDQN ofrece alguna ventaja empírica sobre un clasificador supervisado bien ajustado en condiciones controladas. En segundo lugar, proporciona una comprobación de cordura sobre la dificultad de la tarea en el benchmark, ya que un Random Forest que alcanza puntuaciones muy altas contextualiza lo que QRDQN logra. En tercer lugar, fundamenta cualquier discusión sobre el diseño de la función de recompensa sensible al coste en evidencia concreta. Si ambos modelos se evalúan utilizando las mismas métricas ponderadas por coste, la comparación revela si la función de recompensa de RL conduce realmente al agente hacia un punto de operación diferente al de un modelo supervisado entrenado con los mismos datos, o si ambos enfoques convergen a un comportamiento similar.

\subsection{Evaluación con control de fugas}
El sexto objetivo específico es diseñar pipelines de preprocesamiento y protocolos de validación que impidan estrictamente la fuga de información en todos los niveles de evaluación. Esto requiere la exclusión programática de campos identificativos, como las direcciones IP de origen y destino, las marcas temporales absolutas, los Flow ID y los proxies directos de puertos, en la fase de carga de datos antes de que comience cualquier entrenamiento de modelos \parencite{Arp2020DosDontsMLSecurity,Kaufman2011LeakageDataMining}. Estos campos se excluyen a nivel de esquema mediante el contrato canónico de características, en lugar de depender de un filtrado ad hoc que podría aplicarse de forma inconsistente en distintas ejecuciones o versiones del conjunto de datos.

Las transformaciones de escalado deben ajustarse exclusivamente sobre la partición de entrenamiento y luego aplicarse a las particiones de validación y prueba utilizando únicamente los parámetros aprendidos de los datos de entrenamiento. Ajustar un escalador sobre el conjunto de datos completo antes de la partición permite que la información distribucional del conjunto de prueba influya en la normalización de las características de entrenamiento, una forma sutil de fuga de información que puede inflar el rendimiento reportado en conjuntos de datos con una estructura temporal o de escenario marcada. El escalador ajustado se persiste como parte del artefacto de la ejecución para que pueda recargarse y aplicarse de forma idéntica durante la inferencia de la Fase 2 sobre tráfico de laboratorio, garantizando que se utilice la misma normalización para flujos fuera de distribución sin que ninguna información sobre su distribución llegue al estado del escalador.

\subsection{Análisis de escala de entrenamiento y eficiencia de datos}
El séptimo objetivo específico es evaluar el rendimiento tanto del agente de RL como de las líneas base supervisadas bajo distintas cantidades de datos de entrenamiento, identificando las características de eficiencia de datos de la formulación propuesta y evaluando con qué rapidez converge QRDQN en comparación con las alternativas supervisadas. Los métodos de RL profundo pueden ser muy demandantes en cuanto a muestras, y su dinámica de aprendizaje difiere sustancialmente de la de los métodos de árboles de decisión en conjunto. Comprender cuánto tráfico etiquetado requiere cada enfoque para alcanzar un rendimiento útil es prácticamente relevante: el tráfico de ataque etiquetado es costoso de recopilar, y los escenarios de despliegue reales pueden no tener acceso a los grandes conjuntos de datos equilibrados disponibles en los benchmarks académicos \parencite{Ring2019DatasetSurvey}.

Los experimentos de escala de entrenamiento utilizan subconjuntos de entrenamiento progresivamente más grandes extraídos de la misma partición, midiendo cómo cambian la precisión, el recall, el F1, la tasa de falsos positivos y la tasa de falsos negativos a medida que crece el presupuesto. La comparación entre QRDQN y Random Forest en cada nivel de presupuesto revela si la formulación de RL tiene un requisito mayor de muestras, si converge más lentamente, o si alcanza una meseta de rendimiento final diferente. Estos resultados se interpretan como evidencia interna sobre esta combinación específica de formulación y benchmark, no como afirmaciones generales sobre la eficiencia de datos de RL en todos los contextos de NIDS.

\subsection{Validación complementaria con tráfico de laboratorio}
El octavo objetivo específico es construir un pipeline de inferencia offline robusto capaz de ingerir tráfico capturado por el operador en un entorno de laboratorio doméstico cerrado y no adversarial, y producir predicciones a partir de un modelo entrenado en CICIDS2017. Este nivel de validación ofrece una indicación parcial, de validez externa limitada, de si el esquema canónico de características y el modelo entrenado se comportan de forma razonable sobre flujos extraídos de un entorno de red distinto del benchmark, potencialmente utilizando una versión distinta de CICFlowMeter o una configuración de captura diferente. El cambio de dominio entre la distribución de entrenamiento y la distribución de inferencia es un reto fundamental para los modelos de NIDS y rara vez se evalúa en estudios que solo evalúan sobre un único conjunto de datos público \parencite{Apruzzese2022CrossEvaluationNIDS,Layeghy2023CrossDomainNIDS}.

El pipeline de inferencia aplica recorte opcional por percentiles y z-score para gestionar los valores de características brutas fuera de distribución que caen fuera del rango observado durante el entrenamiento en CICIDS2017. El escalador ajustado de la Fase 1 se carga y aplica sin reajuste, preservando la integridad de la separación entre entrenamiento y prueba. La máscara de presencia/ausencia en el vector de observación codifica explícitamente cualquier característica que no haya podido calcularse a partir de la captura de laboratorio, haciendo visible el caso de datos incompletos en la entrada de predicción en lugar de imputarlo silenciosamente. El resultado esperado de este objetivo es un conjunto de artefactos de predicción y métricas que permitan comparar directamente los resultados del tráfico de laboratorio con los de la validación de la Fase 1, proporcionando evidencia concreta sobre cómo varía el rendimiento bajo un cambio de dominio.

\section{Alcance del trabajo}
El trabajo está delimitado por los requisitos de una evaluación offline de un pipeline de aprendizaje automático tabular. Comprende la ingesta de registros de flujo en formato CSV tanto de CICIDS2017 como de capturas de laboratorio, la aplicación de escalado estándar y recorte opcional de valores atípicos, la ejecución de bucles de entrenamiento mediante el framework Stable Baselines3 (\texttt{sb3-contrib}), y la generación de artefactos de ejecución estructurados y reproducibles que persisten modelos, métricas, escaladores y ficheros de configuración junto a cada ejecución de entrenamiento. La experimentación interna se realiza íntegramente sobre el conjunto de datos CICIDS2017, utilizando particiones explícitas entre entrenamiento y prueba y una escalera de validación multietapa para evaluar el comportamiento del modelo bajo condiciones progresivamente más estrictas antes de probarlo con tráfico capturado externamente.

El alcance incluye explícitamente la escalera de validación multietapa descrita en los objetivos de evaluación. Esta escalera cubre particiones aleatorias estratificadas para comprobaciones básicas de aprendizaje, una comprobación de etiquetas barajadas para verificar que el rendimiento colapsa cuando las etiquetas no contienen señal real, una partición por día o por CSV para comprobar la generalización dentro de CICIDS2017, una evaluación \textit{leave-one-CSV-out} para comprobar la sensibilidad a ficheros de captura específicos, y la inferencia offline sobre flujos capturados por el operador en un laboratorio doméstico cerrado. Cada nivel aporta evidencia diferenciada y forma parte del alcance experimental acordado.

\section{No objetivos}
Esta tesis excluye explícitamente el desarrollo de monitorización de red en línea o la interceptación de paquetes en tiempo real. El sistema no se integrará con cortafuegos, reglas de \texttt{iptables}, controladores de Software-Defined Networking (SDN), ni sistemas de prevención de intrusiones (IPS) activos. Estas exclusiones existen porque el despliegue en producción requeriría un diseño de sistema sustancialmente diferente, una validación de seguridad y una metodología de pruebas operacionales que quedan fuera del alcance de un estudio académico controlado. La tesis está diseñada como un prototipo experimental, no como un producto operativo destinado a su uso en producción.

La tesis tampoco pretende resolver el problema más amplio de las operaciones cibernéticas adversariales multiagente o la interrupción de ataques con dependencia de secuencias, donde el estado del entorno cambia dinámicamente en respuesta a las acciones del defensor. En la formulación actual, cada registro de flujo se clasifica de forma independiente, y el conjunto de datos no reacciona a las decisiones del agente. Esta es una limitación conocida del enfoque \textit{conjunto-de-datos-como-entorno}, e implica que el sistema no puede modelar la adaptación adversarial, la persistencia del atacante ni las cadenas de ataques multietapa que abarcan muchas decisiones individuales de flujo. Por último, el objetivo no es demostrar que QRDQN represente el estado del arte definitivo para todas las tareas de NIDS, sino evaluar su comportamiento bajo un marco experimental cuidadosamente controlado, libre de fugas y reproducible, y compararlo honestamente con líneas base supervisadas.

\section{Resultados esperados}
El resultado esperado de esta tesis es un pipeline de inferencia offline completamente funcional, documentado y reproducible, respaldado por un análisis empírico detallado. El proyecto producirá artefactos persistentes para todas las ejecuciones de entrenamiento y validación, incluyendo ficheros de modelo entrenado, escaladores ajustados, informes de métricas, salidas de validación y registros de configuración que especifican completamente las condiciones bajo las cuales se produjo cada resultado. Esta disciplina en la gestión de artefactos forma parte de la contribución metodológica: un resultado que no puede reproducirse a partir de su rastro de artefactos proporciona una evidencia más débil que uno que sí puede.

En cuanto al contenido empírico, el proyecto producirá mediciones de precisión, recall, F1, tasa de falsos positivos y tasa de falsos negativos tanto para el agente QRDQN como para la línea base de Random Forest en todos los niveles de validación. Los experimentos de escala de entrenamiento producirán curvas de aprendizaje que caracterizarán cómo evoluciona el rendimiento de cada modelo a medida que crece el presupuesto de entrenamiento. Los resultados de inferencia de la Fase 2 caracterizarán cómo varía el rendimiento cuando el modelo se enfrenta a tráfico capturado por el operador en un laboratorio doméstico cerrado y no adversarial, en lugar de flujos de CICIDS2017 reservados para la prueba. En conjunto, estos resultados cuantificarán los límites de generalización de los modelos de RL entrenados en benchmarks cuando se someten a particiones temporales estrictas, evaluación \textit{leave-one-CSV-out} y tráfico de laboratorio doméstico generado por el operador, y proporcionarán una evaluación empírica honesta de si la formulación de RL sensible al coste ofrece ventajas medibles sobre una línea base supervisada bien equiparable.

\chapter{Estado del arte}

\section{Introducción y alcance del capítulo}

Este capítulo mapea el trasfondo técnico de un defensor de ciberseguridad basado en aprendizaje por refuerzo que opera sobre observaciones de red a nivel de flujo. El alcance es deliberadamente más amplio que una comparación de clasificadores binarios. La tesis combina detección de intrusiones en red, representación de características de flujo, conjuntos de datos de referencia públicos, líneas base de aprendizaje supervisado, una formulación de conjunto de datos como entorno al estilo Gymnasium, QRDQN, preprocesamiento con control de fugas, comprobaciones de validación más estrictas, experimentos de eficiencia de datos, y validación planificada o preferida sobre tráfico capturado por el operador en un laboratorio doméstico cerrado y no adversarial.

El capítulo sigue por tanto un embudo. Primero revisa la detección de intrusiones en red y la representación del tráfico, luego sitúa a CICIDS2017 entre los conjuntos de datos públicos, después discute las líneas base de aprendizaje supervisado y aprendizaje profundo, y finalmente enmarca la formulación de RL tanto dentro de la teoría del aprendizaje por refuerzo como en trabajos previos de RL/DRL para detección de intrusiones y defensa cibernética. Las últimas secciones se centran en los riesgos metodológicos, los protocolos de evaluación, la evidencia sobre escala de entrenamiento y la brecha de investigación que aborda esta tesis.

La afirmación central es cautelosa. El aprendizaje por refuerzo para la detección de intrusiones ya existe, y los benchmarks públicos de NIDS son ampliamente utilizados \parencite{Yang2024DRLNIDSSurvey,Kheddar2025RLIDSReview,Gueriani2024DRLIoTIDSSurvey}. Esta tesis no pretende introducir el primer IDS basado en RL, demostrar que QRDQN es el estado del arte para NIDS, ni demostrar una prevención en línea lista para producción. Estudia un pipeline de decisión reproducible, offline, a nivel de flujo, con acciones \texttt{PERMIT}/\texttt{BLOCK}, y evalúa cómo se comporta dicha formulación bajo benchmarks internos controlados, comparación con líneas base y validación progresivamente más estricta.

La revisión también separa tres niveles que con frecuencia se confunden. El primero es la monitorización: recopilar paquetes, flujos, registros o eventos derivados y decidir si son sospechosos. El segundo es la predicción: entrenar un modelo que asigne representaciones de tráfico a etiquetas benignas o maliciosas. El tercero es el control: tomar acciones que modifiquen el estado futuro de la red. Esta tesis trabaja principalmente en los dos primeros niveles. Utiliza el lenguaje de \texttt{PERMIT} y \texttt{BLOCK} porque el entorno expone una decisión de seguridad binaria, pero la implementación actual no es un plano de control en línea. Esa distinción condiciona todas las afirmaciones posteriores.

\section{Detección de intrusiones y monitorización de seguridad en red}

Los sistemas de detección de intrusiones monitorizan la actividad de hosts o redes en busca de evidencias de ataques, usos indebidos o violaciones de políticas \parencite{ScarfoneMell2007}. Un sistema de detección de intrusiones en host observa endpoints, eventos del sistema operativo, registros, procesos, cambios en la integridad de archivos o llamadas al sistema. Un sistema de detección de intrusiones en red observa el tráfico que atraviesa enlaces, taps, puertos de mirror de switches, brokers de paquetes, interfaces de red virtual o colectores de flujos. Ambas perspectivas son complementarias: el HIDS puede ver el contexto local de procesos y archivos que los sensores de red no pueden, mientras que el NIDS puede inspeccionar movimientos laterales, escaneos, tráfico de mando y control, y patrones de comunicación agregados sin instalar un agente en cada host.

También es necesario distinguir el IDS de los sistemas de prevención de intrusiones. Un IDS detecta y notifica principalmente actividad sospechosa. Su salida puede ser una alerta, un evento de registro, un registro de flujo enriquecido, un ticket, un evento SIEM o un artefacto forense. Un IPS se posiciona en línea y puede bloquear, descartar, modificar, limitar la tasa, reiniciar o redirigir el tráfico \parencite{ScarfoneMell2007}. En la práctica, los productos y arquitecturas frecuentemente difuminan esta frontera porque un motor de detección puede alimentar reglas de firewall, acciones EDR, playbooks SOAR o actualizaciones de política SDN. Para la evaluación, sin embargo, la distinción importa. Un detector puede tolerar un perfil de falsos positivos diferente al de un bloqueador en línea, dado que las alertas falsas principalmente afectan a la carga de trabajo del analista, mientras que los bloqueos falsos pueden interrumpir tráfico legítimo de negocio.

Las taxonomías clásicas de IDS distinguen enfoques basados en firmas, basados en anomalías, basados en especificaciones e híbridos \parencite{ScarfoneMell2007}. La detección basada en firmas coteja patrones maliciosos conocidos: secuencias de bytes, predicados de reglas, condiciones de protocolo, cadenas de comandos, rastros de exploits o combinaciones de comportamiento conocido. Su fortaleza es la precisión y la explicabilidad ante amenazas conocidas. Su debilidad es la dependencia del mantenimiento de reglas y la escasa cobertura de ataques nuevos, comportamiento polimórfico, cargas útiles cifradas y variaciones específicas del entorno.

La detección basada en anomalías modela el comportamiento esperado y señala las desviaciones \parencite{SommerPaxson2010ClosedWorld,Ring2019DatasetSurvey}. Esta es la familia más naturalmente asociada con el NIDS basado en ML, pero la palabra anomalía no debe interpretarse como sinónimo automático de aprendizaje no supervisado. Muchos estudios públicos de NIDS son clasificadores supervisados entrenados sobre flujos etiquetados como benignos y de ataque, incluso cuando la motivación más amplia es la detección de anomalías. El reto operativo es que el tráfico benigno es amplio y no estacionario, la prevalencia de ataques es baja, las etiquetas son incompletas, y los cambios en la actividad de negocio pueden parecer anómalos sin ser maliciosos.

La detección basada en especificaciones codifica el comportamiento permitido o esperado, por ejemplo máquinas de estado de protocolo, secuencias de comandos permitidas o flujos de trabajo de control industrial restringidos. Puede ser potente cuando el sistema monitorizado tiene un comportamiento estrecho y estable, pero requiere conocimiento del dominio y mantenimiento. Los sistemas híbridos combinan estas ideas, como reglas de firmas para exploits conocidos, puntuaciones estadísticas de anomalías para comportamiento de flujo inusual, y telemetría de host para confirmación en el endpoint. Los programas modernos de monitorización raramente dependen de un único paradigma de detección.

La generación de alertas y la respuesta activa forman un eje separado. Un sistema puede limitarse a producir alertas para su posterior clasificación, puede enriquecer las alertas con evidencia contextual, puede recomendar pasos de respuesta o puede desencadenar contención activa. Esta tesis se sitúa deliberadamente en el extremo conservador de ese eje. Estudia un modelo de decisión offline similar a un NIDS sobre filas de flujo extraídas. El modelo emite predicciones \texttt{PERMIT}/\texttt{BLOCK} durante los experimentos, pero no descarta paquetes, no reescribe políticas de firewall, no actualiza reglas SDN, no pone en cuarentena hosts ni opera en un bucle de retroalimentación en vivo. La bibliografía de NIDS es, por tanto, relevante, pero las afirmaciones sobre IPS en producción quedan fuera del alcance soportado.

\section{Representación del tráfico de red}

El tráfico de red puede representarse como paquetes, bytes de carga útil, sesiones, conexiones o flujos. Las representaciones a nivel de paquete preservan la temporización, las cabeceras, los tamaños, las direcciones, los flags y, en ocasiones, parte de la carga útil por paquete. Son adecuadas para el análisis de protocolos a grano fino y el modelado de secuencias, pero resultan costosas de almacenar y procesar en redes de alta velocidad. Las representaciones a nivel de carga útil retienen el contenido de la aplicación y pueden identificar exploits o firmas de malware invisibles en los metadatos. Su utilidad está limitada por el cifrado, las restricciones de privacidad, las restricciones legales y el coste de CPU.

Las representaciones a nivel de sesión o conexión agrupan paquetes pertenecientes al mismo intercambio, habitualmente utilizando la quíntupla de IP de origen, IP de destino, puerto de origen, puerto de destino y protocolo. Pueden incluir el estado de la conexión, el comportamiento del handshake, metadatos de la aplicación o atributos derivados de TLS/HTTP. Las representaciones a nivel de flujo son más compactas: agregan paquetes en una ventana temporal o registro similar a una conexión, y resumen el comportamiento mediante duración, bytes, paquetes, direccionalidad, tiempos entre llegadas, estadísticas de longitud de paquetes, conteos de flags TCP y periodos activos o inactivos \parencite{Sperotto2010FlowIDS,Umer2017FlowIDS}.

Los registros de estilo NetFlow e IPFIX son el punto de referencia operativo para la monitorización de flujos escalable. Routers, switches, sondas y colectores exportan registros de flujo compactos que pueden retenerse durante largos periodos y procesarse a escala. Estos registros son atractivos para el NIDS porque reducen el volumen y evitan la inspección de la carga útil. También tienen limitaciones porque los despliegues exportan a menudo solo un subconjunto de los campos posibles, pueden muestrear el tráfico y pueden omitir estadísticas más ricas utilizadas en conjuntos de datos académicos. Esto crea una brecha entre la telemetría de flujos en producción y los CSV de benchmark con abundantes características.

Las características al estilo CICFlowMeter se sitúan en el lado más rico y orientado a la investigación de esa brecha. CICIDS2017 proporciona PCAPs y CSV de flujos bidireccionales generados, con características como la duración del flujo, conteos de paquetes y bytes en sentido directo e inverso, estadísticas de tiempos entre llegadas, estadísticas de longitud de paquetes, conteos de flags TCP, agregados de longitud de cabecera, tasas de flujo y temporización activa/inactiva \parencite{Sharafaldin2018CICIDS2017,Lashkari2017CICFlowMeter,CICIDS2017Web}. Estas características son convenientes para ML tabular porque cada flujo se convierte en una fila de anchura fija, pero dependen de la reconstrucción offline de PCAPs y de las decisiones de extracción de características.

La idoneidad de la representación a nivel de flujo es, por tanto, condicional. Es adecuada para esta tesis porque el defensor implementado consume vectores numéricos de anchura fija, porque CICIDS2017 ya está disponible en formato de flujo, y porque la ruta prevista en Fase 2 también extrae CSV de flujos antes de la inferencia. Está limitada porque se pierden la semántica de la carga útil, el ordenamiento exacto de paquetes, el contenido de la aplicación y algunas dependencias temporales de múltiples flujos. Los ataques cuya evidencia es principalmente a nivel de contenido, a nivel de host o de comportamiento a largo plazo pueden ser solo indirectamente visibles para un clasificador de flujos \parencite{Sperotto2010FlowIDS,Umer2017FlowIDS}.

La disponibilidad de características es otra limitación. Una característica calculada por CICFlowMeter a partir de PCAPs completos puede no estar disponible en un exportador NetFlow/IPFIX en producción, o puede calcularse de manera diferente. Sarhan et al. abordan este problema mapeando conjuntos heterogéneos de características de NIDS hacia una representación estándar alineada con NetFlow para mejorar la generalizabilidad y la explicabilidad \parencite{Sarhan2021}. Esta tesis no adopta el conjunto exacto de características estándar de Sarhan, pero sigue la misma motivación metodológica: definir un contrato de características estable, documentar qué está presente o ausente, y evitar tratar un esquema CSV específico de un benchmark como si fuera universal.

El proyecto actual utiliza un esquema canónico de flujo con 76 valores de características y una máscara de presencia/ausencia de 76 valores, produciendo observaciones de 152 dimensiones. La máscara registra si cada valor canónico estaba presente y era válido, en lugar de ocultar silenciosamente los campos ausentes o imputados. Esto importa al pasar de CICIDS2017 a tráfico de laboratorio porque diferentes pipelines de extracción pueden no producir exactamente las mismas columnas. La máscara no hace que los datos faltantes sean inofensivos, y no resuelve el cambio de dominio. Hace auditable la entrada del modelo: un lector puede distinguir un cero medido de un marcador de posición imputado.

La representación a nivel de flujo también modifica el perfil de privacidad y de fugas. La eliminación de las cargas útiles puede reducir la exposición de privacidad, pero los metadatos aún pueden identificar usuarios, servicios, horarios y escenarios. Las direcciones IP, las marcas temporales absolutas, los Flow ID, la identidad de los endpoints, la estructura del día de captura y los puertos usados como proxies de escenario pueden filtrar etiquetas en lugar de representar comportamiento de ataque \parencite{Arp2020DosDontsMLSecurity,Kaufman2011LeakageDataMining}. La tesis trata por tanto la representación y la selección de características como decisiones de diseño relevantes para la seguridad, no solo como detalles de preprocesamiento.

\section{Conjuntos de datos públicos para la detección de intrusiones en red}

Los conjuntos de datos públicos de NIDS permiten la experimentación compartida, la reproducibilidad y la comparación entre trabajos \parencite{Ring2019DatasetSurvey,Thakkar2020DatasetsReview}. También condicionan las conclusiones. Un conjunto de datos refleja su generador de tráfico, su topología, su proceso de captura, su extractor de características, sus etiquetas, sus scripts de ataque y su formato de distribución. El rendimiento sobre un benchmark público es, por tanto, evidencia sobre una tarea controlada, no evidencia directa de preparación para el despliegue \parencite{Apruzzese2022CrossEvaluationNIDS,Layeghy2023CrossDomainNIDS}.

KDDCup99 y NSL-KDD son históricamente importantes porque dieron forma a la evaluación temprana de IDS \parencite{KDDCup1999,Tavallaee2009NSLKDD}. KDDCup99 contiene registros de conexión derivados del tráfico de la era DARPA, con un esquema de 41 características y categorías de ataque como DoS, Probe, R2L y U2R. NSL-KDD redujo algunos problemas de redundancia y facilitó el uso del benchmark, pero mantuvo el mismo origen básico y lógica de características. Estos conjuntos son útiles para comprender la historia de la evaluación de IDS, pero evidencia débil para la defensa moderna basada en flujos: su tráfico, comportamiento de ataques y esquema son antiguos \parencite{Tavallaee2009NSLKDD,Ring2019DatasetSurvey}. En este repositorio, NSL-KDD es contexto histórico, no la ruta final del defensor actual.

UNSW-NB15 es un conjunto de datos de cyber-range más reciente, generado con tráfico benigno y malicioso contemporáneo, capturas sin procesar y características de ingeniería \parencite{MoustafaSlay2015UNSWNB15,Moustafa2016UNSWEval}. Sigue siendo útil porque amplió el campo más allá de los benchmarks al estilo KDD e incluye categorías de ataque más ricas. CICIDS2017 y CSE-CIC-IDS2018 son conjuntos de datos de la familia CIC ampliamente utilizados para la investigación de NIDS basada en flujos; ambos proporcionan escenarios de ataque generados en laboratorio y características de flujo derivadas de capturas de paquetes \parencite{Sharafaldin2018CICIDS2017,CSECICIDS2018AWS}. Bot-IoT y ToN\_IoT amplían el panorama hacia IoT e IIoT, donde el comportamiento de los dispositivos, la telemetría y los patrones de ataque difieren de las redes empresariales \parencite{Koroniotis2019BotIoT,Moustafa2020ToNIoT,Alsaedi2020ToNIoT}.

\begin{table}[htbp]
\centering
\scriptsize
\setlength{\tabcolsep}{3pt}
\begin{tabular}{p{1.85cm}p{1.9cm}p{2.45cm}p{2.65cm}p{2.8cm}}
\toprule
Conjunto de datos & Papel como benchmark & Representación & Principales fortalezas & Principales advertencias \\
\midrule
KDDCup99 & Benchmark histórico heredado & Registros de conexión con 41 características artesanales & Muy ampliamente utilizado; línea base sencilla para la bibliografía clásica de IDS \parencite{KDDCup1999} & Tráfico obsoleto; redundancia; débil como aproximación al NIDS moderno \parencite{Ring2019DatasetSurvey} \\
NSL-KDD & Benchmark heredado depurado & Mismo esquema de conexión al estilo KDD & Reduce algunos problemas de registros duplicados; útil para comparación histórica \parencite{Tavallaee2009NSLKDD} & Sigue heredando las suposiciones antiguas de DARPA/KDD y el diseño de características \\
UNSW-NB15 & Conjunto de datos moderno de laboratorio/cyber-range & PCAPs, registros y características tabulares de ingeniería & Taxonomía de ataques más contemporánea y características más ricas \parencite{MoustafaSlay2015UNSWNB15} & Entorno controlado; no constituye evidencia directa de producción \\
CICIDS2017 & Benchmark interno principal de la tesis & PCAPs y CSV de flujos CICFlowMeter & PCAPs públicos; múltiples días de ataque; ricas características de flujo \parencite{Sharafaldin2018CICIDS2017} & Problemas conocidos de etiquetado, extracción, duplicación y sensibilidad a la partición \parencite{Engelen2021CICIDSIssues,Lanvin2023Errors} \\
CSE-CIC-IDS2018 & Benchmark de mayor tamaño de la familia CIC & PCAPs y flujos al estilo CICFlowMeter & Colaboración más amplia CIC/CSE; comparación moderna útil \parencite{CSECICIDS2018AWS} & Advertencias similares sobre generación en laboratorio y artefactos de flujo \parencite{Liu2022ErrorPrevalenceNIDS} \\
Bot-IoT & Benchmark de IoT/botnets & PCAPs, flujos al estilo Argus, exportaciones CSV & Escenarios de botnet, escaneo, DoS/DDoS y exfiltración orientados a IoT \parencite{Koroniotis2019BotIoT} & Fuerte desbalanceo y artefactos de topología específicos del conjunto de datos \\
ToN\_IoT & Benchmark de telemetría IoT/IIoT & Flujos de red, registros de host y telemetría de sensores & Perspectiva multimodal de IoT/IIoT \parencite{Alsaedi2020ToNIoT} & Entorno de nicho; los artefactos de identificadores y escenarios siguen siendo posibles \\
\bottomrule
\end{tabular}
\caption[Comparación de conjuntos de datos públicos de NIDS]{Comparación compacta de los conjuntos de datos públicos de NIDS relevantes para la tesis.}
\label{tab:nids-datasets}
\end{table}

La Tabla~\ref{tab:nids-datasets} no establece un rango. Muestra que la elección del conjunto de datos debe alinearse con la solidez de las afirmaciones. Los benchmarks históricos sitúan el campo; los conjuntos de datos modernos de laboratorio permiten comparación controlada; los de IoT prueban suposiciones de tráfico diferentes; y el tráfico de laboratorio externo o similar al de producción aporta evidencia de generalización más sólida. Ninguno de estos conjuntos de datos debe tratarse como un sustituto completo de todas las redes reales. Su valor es máximo con preprocesamiento transparente, líneas base sólidas y protocolos que prueben el cambio de distribución, no solo el rendimiento sobre distribuciones aleatorias emparejadas.

\section{CICIDS2017 como benchmark interno principal}

CICIDS2017 es el benchmark interno principal de esta tesis porque proporciona tráfico benigno y de ataque etiquetado, PCAPs sin procesar y CSV de flujos generados con características al estilo CICFlowMeter \parencite{Sharafaldin2018CICIDS2017,Lashkari2017CICFlowMeter,CICIDS2017Web}. Fue diseñado para mejorar conjuntos de datos de IDS anteriores mediante perfiles de tráfico contemporáneos, actividad benigna realista y múltiples escenarios de ataque \parencite{Sharafaldin2018CICIDS2017,Sharafaldin2019Detailed}. Estas propiedades lo hacen adecuado para la experimentación controlada con modelos de NIDS a nivel de flujo.

Su estructura también es atractiva para la tesis. El conjunto abarca múltiples días de captura y familias de ataque, y está disponible tanto como capturas de paquetes sin procesar como en CSV generados. Los PCAPs están más cerca de la evidencia de medición original. Los CSV facilitan el ML porque CICFlowMeter ya reconstruyó los flujos bidireccionales y calculó características estadísticas. Esa comodidad también es un riesgo: al trabajar solo con CSV resumidos, se puede dejar de comprobar cómo se generaron los flujos, cómo se asignaron las etiquetas y si las filas mantienen la independencia entre las particiones de entrenamiento y prueba.

CICFlowMeter reconstruye flujos de trazas de paquetes mediante claves de endpoint, direccionalidad, lógica de timeout y cálculos de características. Pequeñas diferencias en los parámetros de extracción pueden cambiar filas, duraciones, conteos y etiquetas. Engelen et al. muestran problemas prácticos en la reconstrucción de flujos y la consistencia de CICIDS2017 \parencite{Engelen2021CICIDSIssues}. Lanvin et al. reportan errores que afectan al rendimiento de detección y que correcciones aparentemente pequeñas pueden producir diferencias métricas significativas \parencite{Lanvin2023Errors,FaultyUseCICIDS2017}. Liu et al. analizan la prevalencia de errores en CIC-IDS-2017 y CSE-CIC-IDS-2018, mientras Rosay et al. ofrecen un análisis más amplio y una discusión orientada a la reconstrucción \parencite{Liu2022ErrorPrevalenceNIDS,Rosay2022CICIDS2017Analysis}. La crítica de Dube es especialmente relevante para la disciplina en las afirmaciones, porque cuestiona las interpretaciones habituales de puntuaciones elevadas sobre los datos resumidos de CIC-IDS 2017 \parencite{Dube2024FaultyCICIDS2017}.

Estas críticas no hacen inútil a CICIDS2017; hacen más identificable el uso irresponsable. Los patrones comunes incluyen tratar los CSV pregenerados como verdad incuestionable, usar particiones aleatorias por filas de todos los días, conservar identificadores o marcas temporales que exponen el contexto del escenario, informar solo de la exactitud, ignorar el desbalanceo de clases y no documentar si los datos son originales, corregidos, reconstruidos o curados. Las variantes corregidas o reconstruidas pueden ser valiosas, pero no son intercambiables con la distribución original de CSV. Una tesis debe nombrar qué variante utiliza y evitar mezclar resultados de diferentes variantes como si procedieran de un único conjunto de datos.

El repositorio actual responde a estas preocupaciones mediante una copia curada y rastreada, una copia de referencia sin rastrear y limpieza a nivel de cargador en \texttt{src/load\_cicids2017.py}. El adaptador realiza coerción numérica, tratamiento de valores no válidos, mapeo canónico y exclusión anti-fugas. También ajusta el escalado sobre la partición de entrenamiento y aplica la transformación ajustada a prueba, en vez de ajustar sobre todos los datos. El mapeo canónico en \texttt{src/canonical\_schema.py} excluye IPs, marcas temporales, Flow ID y puertos específicos de las características canónicas.

El código de evaluación admite particiones aleatorias, por CSV/día y validación \textit{leave-one-CSV-out}. El rendimiento con partición aleatoria sirve para comprobar que modelo e implementación aprenden el benchmark interno. Las particiones por día/archivo más difíciles, las comprobaciones de etiquetas barajadas, la validación \textit{leave-one-CSV-out} y la inferencia offline sobre flujos de laboratorio prueban modos de fallo distintos. Una puntuación alta en el escenario más sencillo no basta para una afirmación de despliegue. En esta tesis, CICIDS2017 debe leerse como benchmark interno principal, no como sustituto de la validación con tráfico capturado por el operador en un laboratorio doméstico cerrado y no adversarial.

\section{Aprendizaje supervisado y aprendizaje profundo para NIDS}

La mayor parte del NIDS basado en flujos se formula como aprendizaje supervisado: cada registro de tráfico se mapea a una etiqueta benigna, maliciosa o de familia de ataque \parencite{LiuLang2019IDSSurvey,Ahmad2021}. Los métodos clásicos incluyen árboles de decisión, Random Forests, \textit{SVMs}, \textit{k-NN}, regresión logística, \textit{Naive Bayes} y conjuntos de árboles con \textit{gradient boosting}. Difieren en sus suposiciones, coste operativo y requisitos de memoria, entrenamiento e inferencia. La regresión logística es simple e interpretable y ofrece una referencia lineal, pero está limitada ante interacciones no lineales. \textit{Naive Bayes} es rápido y sirve como línea base débil, aunque sus supuestos de independencia son irrealistas para estadísticas de flujo correlacionadas. k-NN puede capturar vecindades locales, pero escala mal con tablas grandes y es sensible al escalado de características. Las SVMs pueden ser potentes en conjuntos pequeños, pero resultan costosas con tablas de flujos grandes y exigen escalado cuidadoso de características y selección del kernel.

Los modelos basados en árboles son especialmente importantes para NIDS tabular. Los árboles de decisión y Random Forests manejan interacciones no lineales, escalas heterogéneas, variables irrelevantes y distribuciones mixtas \parencite{LiuLang2019IDSSurvey,Apruzzese2022CrossEvaluationNIDS}. Las familias con \textit{gradient boosting}, como \textit{XGBoost}, \textit{LightGBM} y \textit{CatBoost}, son a menudo aprendices tabulares potentes, pero este capítulo no necesita afirmaciones separadas sobre cada implementación si la tesis no las evalúa. El punto metodológico importante es que un clasificador de flujos basado en RL debe compararse con líneas base supervisadas tabulares competentes, no solo con métodos débiles u obsoletos.

Esto es directamente relevante para la tesis. Si cada observación es una fila de flujo etiquetada y la recompensa deriva de la corrección de la clasificación, un clasificador supervisado es la línea base natural. Un agente QRDQN no puede evaluarse únicamente contra su curva de recompensa, pérdida de entrenamiento o una línea base de clase mayoritaria; debe compararse con métodos supervisados sólidos bajo el mismo esquema de características, protocolo de partición, preprocesamiento y métricas. El repositorio ya incluye una línea base de Random Forest alineada con las particiones de QRDQN y artefactos por barrido. La comparación debe limitarse a métricas respaldadas por esos artefactos y al mismo protocolo, en vez de inferir superioridad general.

El aprendizaje profundo amplía el espacio de diseño mediante MLPs, CNNs, RNNs, LSTMs, autoencoders, mecanismos de atención, modelos al estilo transformer y enfoques basados en grafos \parencite{Xiao2021DLIDSSurvey,Asadullah2023MLIDSSurvey,Sayeeduddin2022AINIDS,Nguyen2022BERTFlowNIDS}. Los MLPs tratan las características de flujo como vectores numéricos genéricos; las CNNs se utilizan cuando bytes de paquetes, matrices de características o patrones locales se codifican en forma estructurada. Las RNNs y LSTMs se orientan a secuencias temporales de paquetes, flujos o eventos. Los autoencoders apoyan frecuentemente la detección de anomalías al aprender representaciones comprimidas del comportamiento benigno y señalar errores de reconstrucción. Los modelos al estilo transformer y las redes neuronales en grafo son direcciones más recientes para representaciones de tráfico secuenciales y estructurales.

Estas arquitecturas pueden alcanzar puntuaciones de benchmark muy elevadas, sobre todo en conjuntos públicos, pero las revisiones identifican debilidades recurrentes: preprocesamiento poco claro, particiones aleatorias favorables, análisis insuficiente del desbalanceo, validación externa limitada y reproducibilidad insuficiente \parencite{LiuLang2019IDSSurvey,Ring2019DatasetSurvey,Arp2020DosDontsMLSecurity}. Un modelo profundo puede sobreajustarse a los artefactos del benchmark igual que un modelo superficial. La familia de modelos por sí sola no resuelve fugas, desbalanceo de clases, disponibilidad de características ni cambio de dominio.

La estandarización de características es otra lección importante del aprendizaje supervisado. Diferentes conjuntos de datos y herramientas exponen conjuntos solapados, pero no idénticos, lo que dificulta la comparación entre conjuntos. Sarhan et al. propusieron un mapeo de características alineado con NetFlow para mejorar la generalizabilidad y la explicabilidad en NIDS \parencite{Sarhan2021}. La tesis no adopta ese estándar exacto, pero sí su motivación: un esquema canónico fijo hace auditable la interfaz del modelo y permite comparar las entradas de flujo de CICIDS2017 y de laboratorio a través de un único contrato de preprocesamiento. La máscara de presencia/ausencia extiende esta idea al hacer visibles las características ausentes al modelo y al evaluador.

\section{Fundamentos del aprendizaje por refuerzo}

El aprendizaje por refuerzo estudia cómo un agente selecciona acciones en un entorno para maximizar el retorno \parencite{SuttonBarto2018RL}. Sus conceptos centrales son los estados u observaciones, las acciones, las recompensas, las políticas, las funciones de valor y las transiciones. En un proceso de decisión de Markov, el estado actual resume la información necesaria para la toma de decisiones futura, y las acciones influyen tanto en la recompensa inmediata como en los estados posteriores \parencite{SuttonBarto2018RL}. Esta propiedad secuencial es lo que distingue el RL de la clasificación supervisada ordinaria.

Una política especifica cómo el agente elige acciones a partir de las observaciones. Una función de valor estima el retorno esperado desde un estado, mientras que una función de valor-acción estima el retorno esperado tras tomar una acción particular en un estado. Las recompensas definen el objetivo de la tarea, y el retorno agrega las recompensas futuras, generalmente con descuento. Estas definiciones importan porque cambiar la función de recompensa cambia aquello para lo que el agente está optimizado. En un contexto de seguridad, una recompensa que penaliza más los falsos negativos que los falsos positivos codifica una preferencia operativa concreta; no es un hecho neutral sobre todas las redes.

La distinción entre tareas episódicas y continuas también es relevante. Las tareas episódicas tienen límites y estados terminales, como un episodio de juego o un escenario simulado de respuesta a incidentes. Las tareas continuas modelan la operación continua, como la monitorización continua de la red. Una formulación estática de conjunto de datos como entorno impone a menudo episodios artificiales: un episodio puede ser un número fijo de filas muestreadas, un recorrido por un conjunto de datos o un archivo. Eso es aceptable para la experimentación, pero es más débil que un entorno donde la acción del defensor modifica el estado siguiente.

El Q-learning clásico es un método de diferencia temporal model-free y off-policy que aprende valores de acción sin requerir un modelo de la dinámica del entorno \parencite{WatkinsDayan1992QLearning}. Los métodos model-free aprenden directamente de la interacción o de transiciones registradas, mientras que los métodos basados en modelo aprenden o utilizan un modelo de transición para la planificación. Los métodos on-policy aprenden sobre la política que está generando el comportamiento actual; los métodos off-policy pueden aprender sobre una política objetivo mientras el comportamiento es generado de manera diferente. DQN y QRDQN heredan el enfoque basado en valores y off-policy. La repetición de experiencias, introducida anteriormente como mecanismo para reutilizar experiencias pasadas, también se volvió central en los sistemas de Q-learning profundo posteriores \parencite{Lin1992ReactiveAgents}.

El RL tabular no es apropiado cuando las observaciones son vectores numéricos de alta dimensión. Una tabla de búsqueda sobre todas las posibles combinaciones de características de flujo es inviable, y pequeños cambios numéricos podrían crear efectivamente estados nuevos. La aproximación de funciones es, por tanto, necesaria. El DRL utiliza redes neuronales para aproximar políticas o funciones de valor sobre espacios de observación grandes. En esta tesis, la observación es un vector de 152 dimensiones: 76 valores canónicos de flujo más 76 valores de máscara de presencia/ausencia. Eso justifica un algoritmo de la familia DQN más que un algoritmo de RL tabular, pero no justifica por sí solo el RL frente al aprendizaje supervisado.

La ciberseguridad puede contener problemas de RL secuenciales genuinos. Un defensor puede aislar un host, parchear un servicio, cambiar el enrutamiento, desplegar una contramedida o elegir dónde inspeccionar a continuación, y esas acciones pueden afectar a las observaciones posteriores y a las oportunidades del atacante. Sin embargo, un conjunto de datos estático de flujos etiquetados es un escenario más débil. Si se muestrea una fila de flujo, se clasifica y va seguida de otra fila que no es causalmente afectada por la acción anterior, el problema se parece más a una clasificación sensible al coste o a un proceso de decisión contextual que a una defensa cibernética autónoma completa \parencite{Levine2020OfflineRL,Fujimoto2019BCQ}. La tesis debe mantener esta distinción de forma explícita.

\section{Q-Learning profundo y aprendizaje por refuerzo distribucional}

El Q-Learning Profundo sustituye una función de valor-acción tabular por una red neuronal, haciendo aplicable el RL basado en valores a observaciones de alta dimensión \parencite{Mnih2015DQN}. DQN popularizó dos mecanismos estabilizadores: la memoria de repetición y una red objetivo separada \parencite{Mnih2015DQN}. La memoria de repetición rompe las correlaciones a corto plazo muestreando transiciones pasadas para el entrenamiento y mejora la reutilización de datos. Una red objetivo estabiliza los objetivos de bootstrapping manteniendo el cálculo del objetivo separado de los pesos de la red online, que cambian rápidamente.

Variantes posteriores de DQN abordaron debilidades conocidas. Double DQN reduce el sesgo de sobreestimación desacoplando la selección de acciones de la evaluación del objetivo \parencite{VanHasselt2016DoubleDQN}. Dueling DQN separa la estimación del valor del estado y la ventaja de la acción, lo que puede ayudar cuando muchas acciones tienen valores similares en un estado \parencite{Wang2016DuelingDQN}. Prioritized Experience Replay muestrea transiciones según la señal de aprendizaje en lugar de hacerlo de forma uniforme \parencite{Schaul2016PER}. Rainbow combinó varias mejoras en un único agente de la familia DQN, incluyendo double Q-learning, redes dueling, repetición priorizada, retornos de múltiples pasos, aprendizaje distribucional y exploración con ruido \parencite{Hessel2018Rainbow}.

El aprendizaje por refuerzo distribucional cambia el objeto que se estima. En lugar de estimar únicamente el retorno esperado, aproxima la distribución de los retornos posibles \parencite{Bellemare2017DistributionalRL}. C51 representa la distribución del retorno sobre un soporte categórico fijo. QRDQN usa regresión cuantílica para aproximar los cuantiles del retorno, evitando la misma suposición de soporte fijo \parencite{Dabney2018QRDQN}. QRDQN no es, por tanto, un paradigma separado de DQN; es un miembro distribucional y basado en valores de la familia DQN.

La motivación de seguridad es intuitiva, pero debe formularse con cuidado. Los falsos positivos y los falsos negativos tienen consecuencias asimétricas, y los errores raros de alto coste pueden importar más de lo que sugiere una puntuación media. Una estimación distribucional del valor puede ser útil para estudiar la estructura del retorno bajo recompensas asimétricas. Puede revelar si una decisión tiene una distribución de retornos estrecha o amplia bajo el modelo aprendido. Sin embargo, el RL distribucional no proporciona automáticamente incertidumbre calibrada, seguridad operativa, comportamiento sensible al riesgo ni superioridad sobre el DQN escalar. Esas propiedades requerirían reglas de decisión explícitas, calibración, objetivos de riesgo en la cola o criterios de evaluación más allá de simplemente entrenar un agente distribucional \parencite{Bellemare2017DistributionalRL,Dabney2018QRDQN}.

En esta tesis, QRDQN es una elección algorítmica exploratoria. Es adecuado para un espacio de acciones binario discreto y observaciones de flujo de dimensión moderadamente alta, y está disponible a través de la implementación de SB3-Contrib utilizada por el proyecto \parencite{SB3ContribQRDQNDocs,GymnasiumDocs}. La elección algorítmica es coherente: dos acciones, observaciones numéricas, aprendizaje de valores off-policy, repetición y un objetivo distribucional. La afirmación evidencial sigue siendo empírica. QRDQN no debe describirse como probadamente superior para NIDS a menos que las comparaciones con el mismo protocolo frente a DQN, Random Forest y otras líneas base apoyen esa afirmación.

\section{Aprendizaje por refuerzo para la detección de intrusiones y la defensa cibernética}

El RL y el DRL para la detección de intrusiones ya forman una bibliografía activa \parencite{Yang2024DRLNIDSSurvey,Kheddar2025RLIDSReview,Gueriani2024DRLIoTIDSSurvey,Adawadkar2022CyberSecurityRLSurvey}. Los trabajos previos incluyen clasificadores de intrusiones al estilo DQN, enfoques actor-crítico, formulaciones de entorno adversarial, agentes de selección de características, sistemas orientados a SDN o IoT, y simulaciones de defensa cibernética autónoma. El campo es heterogéneo, por lo que los trabajos no deben tratarse como directamente comparables a menos que su formulación de tarea, conjuntos de datos, acciones, recompensas y protocolos de validación estén alineados.

Una rama está próxima a esta tesis: la detección de intrusiones con conjunto de datos como entorno. Lopez-Martin et al. reformularon la detección supervisada de intrusiones como un problema de RL profundo tratando los registros etiquetados como estados, las predicciones de clase como acciones y la recompensa como función de la corrección \parencite{LopezMartin2020DRLIDS}. Su trabajo posterior extendió una formulación relacionada de RL offline a varios conjuntos de datos, incluyendo CICIDS2017 y UNSW-NB15 \parencite{LopezMartin2021RBFOfflineRL}. Ren et al. utilizaron DRL para la selección de características y la clasificación en CSE-CIC-IDS2018 \parencite{Ren2022IDRDRL}. Estos trabajos muestran que las formulaciones de RL para IDS tienen precedentes.

Otros trabajos de RL orientados a IDS exploran DQN, Double DQN, DQN adversarial, actor-crítico, SAC, DDPG y diseños al estilo Rainbow \parencite{Caminero2019AERLIDS,AEDDQN2020NSLKDD,AESAC2023RLIDS,Alavizadeh2021,Hossain2025DQIDS}. Los detalles varían ampliamente. Algunos trabajos utilizan clasificación binaria, otros etiquetas de ataque multiclase, selección de características, control de umbral o acciones específicas de SDN. Muchos siguen evaluándose sobre conjuntos de datos públicos estáticos. Eso los convierte en antecedentes relevantes para la tesis, pero también limita la comparabilidad directa. Un trabajo que usa NSL-KDD con particiones aleatorias y recompensas de corrección no es evidencia de que un defensor QRDQN offline generalice de CICIDS2017 a tráfico de laboratorio.

Una segunda rama estudia la defensa cibernética genuinamente secuencial. Los enfoques POMDP y de grafo de ataque modelan acciones defensivas bajo incertidumbre \parencite{Hu2020POMDPDefense}. CybORG proporciona un entorno al estilo gym para agentes cibernéticos autónomos \parencite{Standen2021CybORG}. Los trabajos más amplios de defensa cibernética autónoma discuten agentes que seleccionan contramedidas en redes simuladas o emuladas \parencite{Palmer2023_weak,CSET2023AutonomousCyberDefense}. Estos escenarios son conceptualmente diferentes de esta tesis porque las acciones del defensor pueden alterar el estado futuro. Representan una forma más fuerte de problema de RL que la clasificación estática por filas.

\begingroup
\scriptsize
\setlength{\tabcolsep}{2pt}
\renewcommand{\arraystretch}{1.08}
\begin{longtable}{@{}p{1.75cm}p{2.05cm}p{1.95cm}p{3.55cm}p{3.00cm}@{}}
\caption[Trabajos de RL para NIDS]{Síntesis de trabajos de RL/DRL para NIDS y límites de comparabilidad con esta tesis. La relación indicada se infiere de cada fuente y de la formulación QRDQN offline sobre flujos de CICIDS2017 adoptada en este trabajo.}\label{tab:rl-nids-related-work}\\
\toprule
\textbf{Fuente} & \textbf{Formulación de RL} & \textbf{Datos / ámbito} & \textbf{Límite de evaluación} & \textbf{Relación / diferencia con esta tesis} \\
\midrule
\endfirsthead
\multicolumn{5}{@{}l}{\small\itshape Tabla \thetable{} (continuación)} \\
\toprule
\textbf{Fuente} & \textbf{Formulación de RL} & \textbf{Datos / ámbito} & \textbf{Límite de evaluación} & \textbf{Relación / diferencia con esta tesis} \\
\midrule
\endhead
\midrule
\multicolumn{5}{r@{}}{\small\itshape Continúa en la página siguiente} \\
\endfoot
\bottomrule
\endlastfoot
\textcite{NIDSRL2023} & RL; el algoritmo no es verificable desde el registro primario público. & NIDS genérico; datos no verificables. & El registro no expone protocolo ni métricas comparables. & Antecedente temático, no comparación directa con QRDQN/CICIDS2017. \\
\textcite{RLTechniques2023NIDS} & Técnicas de RL; algoritmo no verificable desde el registro primario público. & NIDS genérico; datos no verificables. & El registro no expone protocolo ni métricas comparables. & Antecedente genérico, sin comparabilidad empírica directa. \\
\textcite{Sanusi2023DRLIDS} & DRL para IDS binario. & NSL-KDD; la tarea multiclase se reduce a binaria. & Tesis offline sobre benchmark; no evidencia de despliegue. & Antecedente de IDS binario con DRL, pero datos e implementación distintos. \\
\textcite{Umer2022RLRLIDS} & DQN y gradiente de política con entrenamiento adversarial. & IDS a nivel de paquete y flujo; CICDDoS2019. & Evaluación de benchmark; sin evidencia de despliegue operativo. & Diseño paquete/flujo más rico y datos distintos de QRDQN/CICIDS2017. \\
\textcite{Cevallos2023DRLIDSBP} & Revisión/taxonomía de DRL-IDS, no un algoritmo único. & Literatura de IDS en IoT. & Síntesis bibliográfica, sin benchmark propio. & Contexto de decisiones DRL en IoT, no línea base empírica. \\
\textcite{DDPG2025AttackDetection} & DDPG. & Detección de ataques de red; datos no verificables. & El registro no expone protocolo ni métricas comparables. & Familia de RL distinta; solo antecedente a nivel de título. \\
\textcite{DRLIDSSDN2025} & DRL-IDS con LFTS-RNN y PC-JTFOA. & SDN; NSL-KDD y WPPD. & Particiones de benchmark, no despliegue SDN vivo. & Tarea, modelo y datos distintos; no equivale a una evaluación QRDQN offline. \\
\textcite{HCLRIDS2025IoMT} & CNN-LSTM con DQN y PPO. & IoMT sanitario; CICIoMT2024. & Benchmark especializado; el trabajo deja la detección \textit{zero-day} como trabajo futuro. & Propuesta IoMT específica, no línea base comparable con CICIDS2017. \\
\end{longtable}
\endgroup

La tesis pertenece, por tanto, principalmente a la primera rama: detección offline a nivel de flujo formulada a través de una interfaz de RL. Puede referenciar la defensa cibernética autónoma como contexto más amplio y dirección futura, pero no debe implicar que el sistema actual realice aprendizaje multiagente, adversarial o de control de red en vivo. La formulación más precisa es que la tesis implementa un benchmark de NIDS con conjunto de datos como entorno y una interfaz de acción de seguridad binaria sensible al coste, no un defensor de cyber-range completo.

\section{Clasificación como RL y decisiones de seguridad binarias}

La formulación \texttt{PERMIT}/\texttt{BLOCK} pertenece a la discusión de clasificación como RL. En el entorno actual, cada observación de flujo se presenta al agente, la acción es binaria y la recompensa se calcula a partir de la acción y la etiqueta de verdad. \texttt{PERMIT} corresponde a aceptar el tráfico como benigno, mientras que \texttt{BLOCK} corresponde a identificarlo como ataque. Un falso negativo permite un flujo de ataque; un falso positivo bloquea tráfico benigno.

Este enfoque tiene un beneficio metodológico real: hace explícitas las decisiones de seguridad sensibles al coste. Los errores de IDS no son simétricos en la práctica. Omitir un ataque y bloquear tráfico benigno tienen significados operativos diferentes, y la compensación preferida depende del entorno defendido \parencite{Lee2002CostSensitiveIDS,Cardenas2006IDSFramework}. El diseño de la recompensa puede codificar esta asimetría directamente, lo que es útil para una tesis que estudia reglas de decisión defensivas en lugar de solo la predicción de etiquetas.

La limitación es igualmente importante. Si el conjunto de datos no reacciona a la acción, entonces \texttt{PERMIT}/\texttt{BLOCK} no es una acción de control activa. Es una etiqueta de decisión offline sobre una fila de flujo registrada o extraída. La fila siguiente no cambia porque el agente haya bloqueado o permitido la fila anterior. La formulación es, por tanto, más próxima a una clasificación con recompensa diseñada que a un firewall, un IPS, un controlador SDN o un agente de defensa cibernética autónoma. Esto no es un defecto si se declara con claridad. Solo se convierte en un defecto si el capítulo reclama en exceso el bloqueo en vivo, la preparación para el despliegue o el control secuencial completo.

La función de recompensa también necesita una interpretación cuidadosa. Una recompensa positiva para un verdadero positivo y una penalización mayor para un falso negativo puede expresar la idea de que los ataques omitidos son costosos. Una penalización por falso positivo puede expresar el coste de interrumpir el tráfico benigno o de sobrecargar a los analistas. Una recompensa de cero o pequeño valor para los verdaderos negativos puede mantener el foco en la detección de ataques. Estas son decisiones de diseño, no economía universal de seguridad. Diferentes organizaciones elegirían distintos costes dependiendo de la continuidad del negocio, el modelo de amenaza, la capacidad del analista y la tolerancia a los ataques omitidos.

Dado que la formulación está próxima a la clasificación sensible al coste, las líneas base supervisadas son obligatorias. Un Random Forest sólido, e idealmente otras líneas base tabulares, ayudan a determinar si QRDQN añade valor empírico más allá de las características canónicas, el preprocesamiento y el diseño de la recompensa. Las mismas métricas deben calcularse tanto a partir de las salidas de RL como de las supervisadas. La recompensa de entrenamiento sola no es suficiente, porque una curva de recompensa puede mejorar mientras la precisión, el recall o el comportamiento de falsos positivos sigue siendo operativamente deficiente. La tesis debe interpretar cualquier ventaja de RL como específica del benchmark y el protocolo, a menos que se disponga de evidencia más amplia.

\section{Riesgos metodológicos en NIDS basado en ML/DL/RL}

La bibliografía de NIDS advierte repetidamente contra sobreinterpretar la exactitud de benchmark. Sommer y Paxson señalan que el ML para detección de intrusiones es difícil de validar en entornos realistas: el tráfico benigno es diverso, los ataques raros, las etiquetas difíciles de obtener y las tasas base operativas difieren de las del benchmark \parencite{SommerPaxson2010ClosedWorld}. Arp et al. identifican riesgos afines en ML para seguridad, como fuga de información, sesgo de muestreo, métricas inapropiadas y diseño experimental débil \parencite{Arp2020DosDontsMLSecurity}. La evidencia sobre generalización entre conjuntos y sobre la elección de datos de entrenamiento refuerza que un resultado dentro de un benchmark no basta para inferir el comportamiento en otro entorno \parencite{Layeghy2023CrossDomainNIDS}.

La fuga de información abarca más que poner por error la misma fila en entrenamiento y prueba: incluye cualquier información relacionada con el objetivo que no estaría legítimamente disponible al predecir \parencite{Kaufman2011LeakageDataMining}. En NIDS puede entrar por identificadores de endpoint, marcas temporales absolutas, Flow ID, puertos específicos del escenario, preprocesamiento global ajustado antes de particionar, flujos duplicados o casi duplicados, y particiones que comparten contexto de captura. Un modelo que aprende que una IP, puerto, día o Flow ID concreto corresponde a un ataque puede puntuar bien sin aprender comportamiento malicioso portable.

La fuga temporal y de escenario es especialmente importante. Los conjuntos públicos de NIDS suelen organizarse por día, archivo de captura, ventana de ataque o escenario con script. Una partición aleatoria por filas puede dejar tráfico muy relacionado a ambos lados de la frontera de entrenamiento y prueba. Aunque las filas no sean duplicados exactos, pueden compartir hosts, herramientas de ataque, regímenes de temporización o artefactos de extracción. Esto importa en CICIDS2017 porque trabajos posteriores reportan errores o artefactos en extracción de flujos, etiquetado y uso de CSV resumidos \parencite{Engelen2021CICIDSIssues,Lanvin2023Errors,Liu2022ErrorPrevalenceNIDS,Rosay2022CICIDS2017Analysis,Dube2024FaultyCICIDS2017}.

El preprocesamiento también puede introducir fugas. El escalado, la imputación, la selección de características, el recorte de valores atípicos o la reducción de dimensionalidad deben ajustarse sobre entrenamiento y aplicarse después a los datos de validación o de prueba. Si un escalador se ajusta sobre todos los datos antes de particionar, la información distribucional de prueba entra en el entrenamiento. Puede parecer menos grave que filtrar etiquetas, pero con estructura fuerte por día o escenario también puede inflar la confianza. Por ello, los objetos \texttt{StandardScaler} ajustados sobre entrenamiento en el repositorio forman parte del diseño metodológico, no solo de la codificación.

El desbalanceo de clases y el problema de la tasa base complican la interpretación. En redes operativas, los ataques suelen ser raros frente al tráfico benigno. Un detector puede lograr exactitud elevada mientras omite muchos ataques, o alto recall mientras genera demasiadas alertas falsas para investigarlas. El argumento de la tasa base de Axelsson sigue siendo relevante: incluso tasas bajas de falsos positivos pueden dominar los flujos de alertas cuando la prevalencia real de ataque es baja \parencite{Axelsson1999BaseRate}. Por eso la exactitud no debe ser la métrica principal para NIDS.

El RL introduce riesgos adicionales: el entrenamiento profundo puede ser estocástico y sensible a semillas aleatorias, hiperparámetros, modelado de la recompensa, configuración del buffer de repetición, exploración y detalles del entorno. En un conjunto de datos estático usado como entorno, buenos resultados pueden reflejar las mismas ventajas de características y partición que benefician a modelos supervisados. Una afirmación de superioridad de QRDQN requiere el mismo protocolo de comparación y reconocer la varianza entre ejecuciones cuando no hay semillas repetidas. Las ganancias de una única ejecución deben tratarse como preliminares salvo que las sostengan experimentos repetidos o artefactos de validación sólidos.

\section{Protocolos de evaluación, métricas y reproducibilidad}

La evaluación debe ajustarse a la solidez de la afirmación. Las particiones aleatorias sobre el mismo conjunto de datos son aceptables para comprobaciones internas de aprendizaje, pero son evidencia débil para el despliegue. Una escalera de validación práctica para esta tesis comienza con particiones aleatorias estratificadas, avanza hacia particiones temporales o por CSV/día, la validación leave-one-scenario o leave-one-CSV-out, las pruebas entre conjuntos de datos con un esquema de características armonizado, y finalmente el tráfico capturado de forma independiente o similar al de producción. La tesis no alcanza este último peldaño; su inferencia sobre tráfico de laboratorio, generado por el operador en un entorno doméstico cerrado y no adversarial, se sitúa por debajo. Cada paso aumenta la separación distribucional entre datos de entrenamiento y prueba, y por tanto permite afirmaciones de generalización más sólidas.

El repositorio actual refleja esta escalera. El entrenamiento con partición aleatoria verifica que la implementación de QRDQN, el esquema de características, la función de recompensa y el bucle de entrenamiento pueden aprender el benchmark controlado. La comprobación B con etiquetas barajadas comprueba si el rendimiento colapsa hacia una línea base de clase mayoritaria cuando las etiquetas ya no contienen señal real. La comprobación C utiliza una partición por CSV/día más difícil. El script leave-one-CSV-out comprueba si retener un CSV original de CICIDS2017 modifica el comportamiento. La inferencia offline de Fase 2 sobre CSV de flujos de laboratorio prueba la cuestión separada de si un modelo entrenado sobre el benchmark puede procesar flujos capturados por el operador en un laboratorio doméstico cerrado.

Las métricas también deben interpretarse operativamente. La exactitud es útil pero insuficiente para NIDS desequilibrados. La precisión indica cuán fiables son las alertas positivas. El recall indica qué proporción del tráfico de ataque es detectada. F1 proporciona un equilibrio compacto, pero oculta la compensación real entre falsos positivos y falsos negativos. La tasa de falsos positivos mapea al tráfico benigno bloqueado o alertado incorrectamente. La tasa de falsos negativos mapea al tráfico malicioso omitido por el detector \parencite{Axelsson1999BaseRate,Fawcett2006ROC,Saito2015PRROC}. Las matrices de confusión, la precisión, el recall, F1, TP, FP, FN, TN, FPR y FNR por clase deben preferirse sobre un único número resumen \parencite{Milenkoski2015IDSEvaluationSurvey}.

La evaluación sensible al coste es relevante porque el punto de operación correcto depende del contexto. La guía del NIST y la bibliografía de evaluación de IDS enfatizan ambas las compensaciones entre falsos positivos y falsos negativos \parencite{ScarfoneMell2007,Cardenas2006IDSFramework}. En esta tesis, el diseño de recompensa asimétrico debe tratarse como una suposición de escenario y evaluarse empíricamente, no como un modelo de coste universal para todas las redes. Si la recompensa penaliza fuertemente los falsos negativos, la política resultante puede preferir bloquear más flujos. Ese comportamiento debe juzgarse mediante precisión, recall, tasa de falsos positivos y las suposiciones de coste específicas declaradas en el experimento.

La reproducibilidad es parte de la evidencia, no un detalle administrativo. Los resultados de ML en ciberseguridad deben preservar la versión del conjunto de datos, la variante exacta del conjunto de datos, el esquema de características, los campos excluidos, el orden del preprocesamiento, la lógica de partición, el estado del escalador, las semillas aleatorias, la configuración del modelo, la configuración de la recompensa, el identificador de ejecución, las métricas y las predicciones donde sea factible. Olszewski et al. muestran que la reproducibilidad sigue siendo un problema importante en la investigación de ML en seguridad \parencite{Olszewski2023ReproSecurityML}. El uso en el repositorio de configuraciones guardadas, métricas, artefactos del modelo, salidas de validación y carpetas de ejecución es, por tanto, una fortaleza metodológica, pero los artefactos faltantes deben seguir claramente marcados como pendientes.

Esta disciplina en los informes también protege contra las afirmaciones excesivas accidentales. Un resultado vinculado a la ejecución C03 de QRDQN con partición aleatoria completa es un resultado para esa ejecución, con esa partición, configuración de recompensa, ruta de preprocesamiento y estado de artefactos. No debe convertirse silenciosamente en una afirmación sobre todas las configuraciones de QRDQN, todas las variantes de CICIDS2017 o todo el tráfico de red. Cuanto más sólida sea la afirmación, más completo debe ser el rastro de artefactos.
% Esto de arriba sobre artefactos faltantes habrá que quitarlo, nunca habrá artefactos faltantes, ya que si faltan, es que no están dentro del alcance de la tesis, y no hay que decirlo. La tesis debe ser clara sobre lo que está dentro del alcance y lo que no. De todos modos, esos artefactos se generarán pronto por el usuario ejecutando los nuevos experimentos.

\section{Eficiencia de datos y evaluación a escala de entrenamiento}

La eficiencia de datos pregunta cuánto tráfico etiquetado se necesita para alcanzar un rendimiento útil, y cómo cambia la curva de aprendizaje a medida que crece el conjunto de entrenamiento. Esta pregunta es importante para NIDS porque el tráfico de ataque etiquetado es costoso, las distribuciones de tráfico cambian y la abundancia de benchmarks públicos no implica abundancia de datos de despliegue \parencite{Ring2019DatasetSurvey,Apruzzese2022CrossEvaluationNIDS}. También es importante para el RL porque los métodos de RL profundo pueden ser exigentes en muestras, especialmente cuando las señales de recompensa son escasas, ruidosas, retardadas o altamente diseñadas \parencite{SuttonBarto2018RL,Hessel2018Rainbow}.

La bibliografía existente de RL-IDS con frecuencia entrena sobre conjuntos de datos públicos completos e informa de puntuaciones finales en el benchmark, mientras que las ablaciones controladas a escala de entrenamiento son menos comunes \parencite{Yang2024DRLNIDSSurvey,Kheddar2025RLIDSReview}. El trabajo de NIDS supervisado con pocos ejemplos y aprendizaje incremental por clases aborda preguntas relacionadas desde una perspectiva diferente \parencite{DiMonda2024FewShotNIDS}, pero la pregunta de la tesis es más estrecha: bajo un mismo esquema de características y protocolo de evaluación, ¿cómo escala QRDQN con las filas de entrenamiento disponibles, y cómo se compara eso con las líneas base supervisadas?

La respuesta importa porque la elección del modelo y las necesidades de datos están conectadas. Si un Random Forest alcanza un rendimiento estable con muchas menos filas que QRDQN bajo la misma partición, eso es evidencia importante incluso si QRDQN lo alcanza posteriormente. Si QRDQN es más sensible a los datos reducidos o al desbalanceo de la clase de ataque, la tesis debe reportarlo como una limitación. Si QRDQN es competitivo con presupuestos menores, la evidencia debe seguir vinculada al benchmark y al protocolo de validación, no generalizada a todos los escenarios de NIDS.

La tesis posiciona los experimentos a escala de entrenamiento como evidencia metodológica. Presupuestos como 100k, 250k, 500k, 1M y 2M muestras pueden comprobar si la formulación de RL requiere sustancialmente más datos que una línea base supervisada para alcanzar un rendimiento comparable. Tales experimentos no deben enmarcarse como prueba de aprendizaje con pocos ejemplos ni de preparación para el despliegue. Son evidencia interna sobre los requisitos de datos de esta formulación particular. La salida útil es una comparación de curvas de aprendizaje con preprocesamiento consistente, semillas aleatorias donde sea factible, métricas compartidas y rutas de artefactos explícitas.

\section{Brecha de investigación y posicionamiento de esta tesis}

La bibliografía respalda una brecha de investigación estrecha pero significativa. El NIDS y el ML basado en flujos son áreas maduras; conjuntos de datos públicos como CICIDS2017 son útiles pero frágiles; las líneas base tabulares supervisadas son sólidas; el RL y el DRL para IDS ya existen; y la defensa cibernética autónoma es un campo más amplio y más secuencial que la implementación actual \parencite{Ring2019DatasetSurvey,LiuLang2019IDSSurvey,Yang2024DRLNIDSSurvey,Kheddar2025RLIDSReview}. La brecha no es la ausencia de RL para la detección de intrusiones.

La contribución de esta tesis es un estudio experimental reproducible de una formulación offline de decisión de seguridad binaria a nivel de flujo. El sistema mapea CICIDS2017 y entradas posteriores de flujos de laboratorio en un esquema canónico fijo de 76 características, aumenta las observaciones con una máscara de presencia/ausencia para producir entradas de 152 dimensiones, expone las filas a través de un entorno compatible con Gymnasium, y entrena un agente QRDQN para elegir entre \texttt{PERMIT} y \texttt{BLOCK}. Esto hace explícita la interfaz de decisión sensible al coste mientras preserva una obligación de comparación clara contra las líneas base supervisadas.

La tesis también está posicionada metodológicamente. Trata a CICIDS2017 como el benchmark interno principal, no como prueba de preparación para la producción. Separa los resultados con partición aleatoria de la validación más estricta. Utiliza preprocesamiento con control de fugas y comprobaciones de validación para reducir el aprendizaje evidente de artefactos. Trata el tráfico de laboratorio capturado externamente como el siguiente paso evidencial preferido, reconociendo al mismo tiempo que la inferencia actual de Fase 2 es offline y no realiza bloqueo activo.

La brecha final es, por tanto, una brecha de integración: un benchmark QRDQN reproducible, con control de fugas y sensible al coste para el NIDS a nivel de flujo, evaluado contra líneas base supervisadas e interpretado a través de una escalera de validación. La tesis puede contribuir haciendo explícito el contrato de características, guardando artefactos de ejecución, documentando las suposiciones de la recompensa, comparando contra Random Forest bajo particiones emparejadas, y mostrando cómo el rendimiento interno de CICIDS2017 cambia bajo una validación más estricta y restricciones de escala de entrenamiento.

La afirmación final defendible no es que QRDQN sea universalmente mejor que los modelos de NIDS supervisados, ni que el proyecto implemente defensa cibernética autónoma. La afirmación defendible es que la tesis evalúa un defensor a nivel de flujo basado en QRDQN, sensible al coste, bajo un benchmark reproducible y un pipeline de validación, y utiliza líneas base supervisadas, comprobaciones con control de fugas, análisis de escala de entrenamiento e inferencia sobre tráfico de laboratorio para determinar hasta qué punto puede confiarse en esa formulación. Si experimentos posteriores muestran un rendimiento débil entre archivos o con tráfico de laboratorio, eso no invalida la tesis. Agudiza su contribución al demostrar los límites del RL entrenado sobre benchmark bajo un cambio de distribución más realista.

\chapter{Metodología}

\section{Visión general metodológica}
Este capítulo describe el diseño metodológico empleado para construir y evaluar el defensor de ciberseguridad basado en aprendizaje por refuerzo (RL) propuesto. El objetivo no es presentar un sistema de prevención de intrusiones (IPS) listo para producción, sino definir un marco experimental controlado, reproducible y offline para decisiones de seguridad binarias sobre registros de flujos de red. Cada flujo se trata como un único punto de decisión: el defensor observa una representación numérica del flujo y decide si permitirlo (\texttt{PERMIT}) como tráfico benigno o bloquearlo (\texttt{BLOCK}) como tráfico malicioso.

La metodología es deliberadamente conservadora. Los benchmarks públicos de detección de intrusiones pueden producir resultados inflados cuando el preprocesamiento, la construcción de las particiones, los controles de fuga de información y las comparaciones con líneas base son débiles \parencite{Arp2020DosDontsMLSecurity,Engelen2021CICIDSIssues,Lanvin2023Errors}. Por ello, el pipeline se ha diseñado en torno a contratos de datos explícitos, artefactos de ejecución reproducibles, normalización ajustada únicamente sobre datos de entrenamiento, y una validación progresivamente más estricta. La misma representación canónica se utiliza tanto para el benchmark interno de CICIDS2017 como para la inferencia offline sobre tráfico de laboratorio en la Fase 2, de modo que la pregunta experimental no queda ligada a una convención de nomenclatura de CSV concreta ni a un extractor de características particular.

A grandes rasgos, la metodología comprende seis etapas encadenadas:
\begin{enumerate}
    \item ingerir los datos CSV de CICIDS2017 a nivel de flujo y mapear las etiquetas a una tarea binaria de defensa;
    \item eliminar los identificadores susceptibles de provocar fuga de información y normalizar los valores numéricos malformados;
    \item mapear las columnas específicas de cada extractor a un esquema canónico de características fijo de 76 características;
    \item añadir una máscara de presencia/ausencia de 76 dimensiones, produciendo observaciones de 152 dimensiones;
    \item entrenar y evaluar un defensor QRDQN compatible con Gymnasium bajo protocolos de partición controlados;
    \item comprobar la robustez mediante comprobaciones de validación, líneas base supervisadas e inferencia offline en la Fase 2.
\end{enumerate}

Este diseño mantiene la tesis dentro de un marco experimental offline. Las salidas \texttt{PERMIT} y \texttt{BLOCK} son predicciones del modelo sobre registros de flujo almacenados. No son órdenes de cortafuegos y no modifican ninguna ruta de red activa.

\section{Fases experimentales}
La metodología experimental se divide en dos fases. La primera fase comprueba el aprendizaje y la generalización interna sobre un benchmark público. La segunda fase comprueba si el pipeline de preprocesamiento y modelo resultante es capaz de procesar tráfico capturado de forma independiente bajo cambio de dominio.

\subsection{Fase 1: benchmark interno de CICIDS2017}
La Fase 1 utiliza CICIDS2017 como benchmark interno principal \parencite{Sharafaldin2018CICIDS2017,CICIDS2017Web}. CICIDS2017 es adecuado para esta tesis porque proporciona tráfico benigno y de ataque etiquetado, capturas de paquetes en bruto y CSVs de flujos derivados de una extracción al estilo CICFlowMeter \parencite{Lashkari2017CICFlowMeter}. También contiene múltiples escenarios de ataque en diferentes días de captura, lo que permite diseños de partición más estrictos que una única partición aleatoria entrenamiento-prueba.

La fase no se trata como un sustituto limpio de un entorno de producción. CICIDS2017 tiene problemas conocidos que incluyen flujos duplicados o inconsistentes, errores de etiquetado y de extracción, y sensibilidad a la partición \parencite{Engelen2021CICIDSIssues,Lanvin2023Errors,Liu2022ErrorPrevalenceNIDS}. Por esa razón, la Fase 1 se enmarca como un benchmark interno con controles explícitos. El pipeline elimina los campos de tipo identificador, ajusta la normalización únicamente sobre los datos de entrenamiento, registra las configuraciones de ejecución y evalúa el modelo mediante una escalera de validación. Una puntuación elevada en la partición aleatoria se interpreta como evidencia de que el modelo es capaz de aprender la tarea del benchmark; no se interpreta como preparación para el despliegue.

\subsection{Fase 2: inferencia offline sobre tráfico de laboratorio}
La Fase 2 evalúa la transferencia más allá de la distribución del benchmark. El tráfico se captura en un entorno de laboratorio privado, se convierte offline en CSVs de flujos, se mapea al mismo esquema canónico, se normaliza con los artefactos de entrenamiento persistidos y se procesa con el modelo entrenado. Esta fase es una prueba de inferencia sobre una distribución externa, no un despliegue en vivo.

El propósito metodológico de la Fase 2 es exponer el cambio de dominio. Los estudios públicos sobre NIDS suelen reportar resultados sólidos dentro del mismo conjunto de datos que se degradan cuando el modelo se evalúa sobre tráfico producido por redes, herramientas de captura, periodos temporales o definiciones de características diferentes \parencite{Apruzzese2022CrossEvaluationNIDS,Layeghy2023CrossDomainNIDS}. La Fase 2 se centra por tanto en la inferencia offline reproducible, los diagnósticos y el seguimiento de artefactos. No garantiza un rendimiento estable en redes empresariales arbitrarias y no implementa bloqueo activo en tiempo real.

\section{Fuente de datos y representación de la entrada}
La tesis utiliza una representación tabular a nivel de flujo. Esta representación es un compromiso deliberado entre relevancia operativa, manejabilidad y privacidad. Captura el comportamiento estadístico a lo largo de secuencias de paquetes evitando la inspección de la carga útil.

\subsection{CSVs de flujos de CICIDS2017}
La fuente de datos principal es la publicación de CSVs de flujos de CICIDS2017. Cada fila de un CSV corresponde a un flujo de red bidireccional resumido mediante estadísticas numéricas como duración, conteos de paquetes, totales de bytes, tiempos de llegada entre paquetes, estadísticas de longitud de paquete y conteos de flags TCP. Las representaciones basadas en flujos se utilizan ampliamente en la investigación sobre NIDS porque comprimen el comportamiento de la red en filas numéricas de anchura fija que pueden procesarse mediante métodos estándar de aprendizaje automático \parencite{Sperotto2010FlowIDS,Umer2017FlowIDS}.

El pipeline de carga local espera los ficheros CSV oficiales de CICIDS2017 y admite presupuestos de filas tanto completos como limitados. Los presupuestos limitados permiten pruebas rápidas de humo e iteración ágil; los presupuestos completos son los utilizados en las ejecuciones finales del benchmark. El cargador lee los CSVs grandes en fragmentos, estandariza los nombres de columna eliminando espacios en blanco inconsistentes, identifica la columna de etiqueta de forma robusta y aplica después la limpieza de características antes de cualquier escalado dependiente de la partición.

\subsection{Mapeo de etiquetas binarias}
CICIDS2017 incluye varias etiquetas de ataque. La tesis las colapsa en una decisión binaria del defensor:
\begin{itemize}
    \item etiqueta \texttt{0}: \texttt{BENIGN}, mapeada a la acción deseada \texttt{PERMIT};
    \item etiqueta \texttt{1}: cualquier clase de ataque, mapeada a la acción deseada \texttt{BLOCK}.
\end{itemize}

Este mapeo binario se corresponde con el modelo de decisión de primera línea estudiado aquí. Además, hace que el espacio de acciones de RL sea sencillo y auditable. La limitación es que la información sobre la familia de ataque no es el objetivo de aprendizaje principal en la ruta de RL actual. Cualquier análisis de errores por familia debe por tanto estar respaldado por etiquetas de familia preservadas o artefactos independientes; no puede inferirse a partir de las etiquetas binarias por sí solas.

\subsection{Representación tabular a nivel de flujo}
El sustrato de aprendizaje es una tabla a nivel de flujo en lugar de paquetes en bruto o bytes de carga útil. Esto reduce la dimensionalidad y hace factible el entrenamiento offline a gran escala. También evita algunas complicaciones legales y técnicas de la inspección profunda de paquetes, especialmente cuando las cargas útiles están cifradas \parencite{Sperotto2010FlowIDS,Umer2017FlowIDS}. Sin embargo, la abstracción de flujo pierde el contenido, el orden preciso de los paquetes y parte del contexto de largo alcance entre múltiples flujos. La metodología estudia por tanto un defensor basado en flujos, no un sistema completo de análisis forense de red.

\section{Limpieza de datos y preprocesamiento}
El preprocesamiento se trata como parte del método experimental, no como un detalle de implementación. En los benchmarks de NIDS, pequeñas decisiones de preprocesamiento pueden cambiar el significado de la tarea al filtrar etiquetas, suprimir flujos malformados o hacer que los datos de prueba influyan en el entrenamiento.

\subsection{Normalización de columnas y coerción numérica}
Las columnas CSV en bruto pueden contener espacios iniciales, capitalización inconsistente y tipos de datos mixtos. El cargador primero normaliza los nombres de columna y después convierte las columnas de características a valores numéricos. La matriz de características final se representa como \texttt{float32}, mientras que las etiquetas se representan como \texttt{int64}. Este contrato de tipos mantiene estable la interfaz del entorno, el escalador y el modelo en las etapas posteriores.

No se permite que las filas introduzcan cadenas de texto arbitrarias en el modelo. Las columnas de características no numéricas que sobreviven a la eliminación de identificadores quedan excluidas. Esto evita la introducción accidental de cadenas categóricas cuya codificación sería no controlada o específica de un conjunto de datos.

\subsection{Valores inválidos, valores ausentes e imputación}
La extracción de flujos de red puede producir valores numéricos inválidos. Las características de tasa, por ejemplo, pueden volverse infinitas o indefinidas cuando se calculan sobre una duración nula o casi nula. El pipeline sustituye los infinitos y los valores numéricos ausentes por cero antes del mapeo canónico. Esta elección es pragmática: muchos contadores de flujo, tasas y estadísticas agregadas pueden utilizar de forma segura el cero como marcador de posición neutro cuando se acompaña de una señal explícita de presencia/ausencia.

La imputación por cero sola sería insegura porque el modelo no podría distinguir un cero medido real de un marcador sintético. El esquema canónico empareja por tanto la imputación con una máscara de presencia/ausencia. La máscara hace visible la disponibilidad de características tanto para el modelo como para los diagnósticos posteriores.

\subsection{Exclusión de campos susceptibles de fuga}
El pipeline de preprocesamiento aplica una política anti-fuga fundamentada en los riesgos de evaluación conocidos en seguridad con ML \parencite{Arp2020DosDontsMLSecurity,Kaufman2011LeakageDataMining}. Se excluyen las direcciones IP de origen y destino, los Flow IDs, las marcas temporales absolutas y los identificadores de endpoints. Los campos de puerto directo también se excluyen porque pueden actuar como indicadores del escenario en conjuntos de datos controlados por scripts. Por ejemplo, un día de ataque puede concentrar tráfico malicioso en servicios concretos, lo que permitiría al modelo memorizar un script de captura en lugar de aprender estadísticas de comportamiento de flujo.

Esta exclusión se aplica antes del mapeo canónico. El modelo recibe por tanto únicamente características estadísticas de flujo y los valores de máscara correspondientes. Si en el futuro se introducen nuevos conjuntos de datos o extractores, sus adaptadores deben preservar la misma política: sin identificadores, sin campos de tiempo absoluto y sin puertos utilizados directamente como atajos de etiqueta.

\section{Esquema canónico de características}
El esquema canónico es el contrato de datos central de la tesis. Separa la interfaz del modelo de los nombres de columna, el orden y la disponibilidad de características de un extractor concreto.

\subsection{Características canónicas de flujo}
El esquema define 76 características numéricas de flujo ordenadas. Incluyen duración, conteos de paquetes hacia adelante y hacia atrás, totales de bytes, estadísticas de longitud de paquete, estadísticas de tiempo entre llegadas de flujo y por dirección, tasas de paquetes y bytes, conteos de flags TCP, longitudes de cabecera, tamaños de ventana, estadísticas de subflujos y características de temporización activa/inactiva. Estas características se han elegido porque son estadísticas de flujo y no identificadores, y porque son ampliamente compatibles con la extracción al estilo CICFlowMeter.

Para CICIDS2017, un adaptador mapea las columnas originales del CSV a estos nombres canónicos. Para la Fase 2, el script de inferencia robusta mapea las columnas del extractor de flujos de laboratorio a los mismos nombres canónicos. El modelo se entrena y evalúa por tanto frente a un contrato de características fijo. Este diseño sigue la motivación metodológica más amplia de la estandarización de características en NIDS: la comparación entre dominios resulta más auditable cuando las definiciones de características son explícitas y estables \parencite{Sarhan2021}.

La Tabla~\ref{tab:canonical-groups} enumera los grupos lógicos dentro del esquema de 76 características y el número de características aportadas por cada grupo. La agrupación es informativa; el modelo recibe los 76 valores como un vector plano sin señales de agrupación estructural.

\begin{table}[htbp]
\centering
\begin{tabular}{lcp{7cm}}
\toprule
Grupo de características & Cantidad & Estadísticas representativas \\
\midrule
Duración del flujo & 1 & Duración total del flujo bidireccional \\
Estadísticas de paquetes hacia adelante & 12 & Conteo de paquetes, bytes totales, media/desv.\ típica/mínimo/máximo de longitudes de paquete \\
Estadísticas de paquetes hacia atrás & 12 & Las mismas estadísticas para la dirección inversa \\
Agregados bidireccionales & 8 & Conteos combinados, bytes totales, promedios de tasa de ráfaga \\
Estadísticas de tiempo entre llegadas & 10 & Media, desv.\ típica, mínimo y máximo del tiempo entre llegadas del flujo, hacia adelante y hacia atrás \\
Conteos de flags TCP & 8 & Conteos de FIN, SYN, RST, PSH, ACK, URG, CWE, ECE \\
Agregados de longitud de cabecera & 4 & Totales de bytes de cabecera hacia adelante y hacia atrás \\
Estadísticas de tamaño de ventana & 3 & Tamaño de ventana inicial y medio por dirección \\
Características de tasa de flujo & 6 & Bytes/s, paquetes/s por dirección y combinados \\
Estadísticas de subflujos & 6 & Promedios de paquetes y bytes de subflujo por dirección \\
Temporización activa/inactiva & 6 & Media, desv.\ típica y máximo de las duraciones de los periodos activos e inactivos \\
\bottomrule
\end{tabular}
\caption{Agrupación lógica de las 76 características canónicas de flujo. Todos los grupos se fusionan en un único vector plano antes de la entrada al modelo.}
\label{tab:canonical-groups}
\end{table}

\subsection{Máscara de presencia/ausencia}
Diferentes extractores pueden no proporcionar todas las características canónicas, e incluso una columna fuente presente puede contener valores inválidos. El pipeline produce por tanto una máscara de presencia/ausencia de 76 dimensiones. Para cada característica canónica $x_i$, el valor de máscara correspondiente $m_i$ es:
\begin{itemize}
    \item $m_i = 1$ cuando la característica fuente está presente y es válida;
    \item $m_i = 0$ cuando la característica está ausente o ha tenido que ser imputada.
\end{itemize}

La máscara desempeña dos funciones metodológicas. En primer lugar, evita la imputación silenciosa: el modelo puede condicionar su decisión a si un valor fue observado o rellenado. En segundo lugar, facilita los diagnósticos de la Fase 2, porque un extractor de laboratorio que carece de muchas características al estilo CICIDS2017 puede detectarse mediante un patrón de máscara diferente en lugar de quedar oculto dentro de un vector numérico denso.

\subsection{Vector de observación final}
La observación final es:
\[
    o = [x_1, x_2, \dots, x_{76}, m_1, m_2, \dots, m_{76}]
\]
Esto produce un vector \texttt{float32} de 152 dimensiones. La primera mitad contiene los valores de las características, y la segunda mitad contiene la máscara. El tamaño de la observación es por tanto invariante entre CICIDS2017, las particiones de validación y la inferencia en la Fase 2. Esta invariancia es importante porque la red de política del QRDQN espera una dimensión de entrada fija.

\section{Separación entrenamiento-prueba y escalado}
La metodología separa el ajuste y la evaluación en cada etapa de preprocesamiento que pueda aprender de los datos. Esto es necesario porque la fuga de información puede producirse no solo a través de columnas, sino también a través de estadísticas como medias globales, varianzas, percentiles y distribuciones de clases.

\subsection{Partición aleatoria estratificada}
La partición aleatoria estratificada es el primer peldaño de validación. Preserva las proporciones de clases entre entrenamiento y prueba y es útil como comprobación de la capacidad de aprendizaje. Si el modelo no puede aprender bajo esta partición favorable, las evaluaciones más estrictas posteriores tienen pocas probabilidades de tener éxito.

Sin embargo, la partición aleatoria por filas es la evidencia más débil de este proyecto. Los conjuntos de datos de flujos pueden contener registros correlacionados, patrones repetidos y artefactos del contexto de captura. Situar flujos adyacentes o muy similares a ambos lados de la partición puede sobreestimar la generalización \parencite{Lanvin2023Errors,Arp2020DosDontsMLSecurity}. Los resultados de la partición aleatoria se reportan por tanto únicamente como evidencia del benchmark interno.

\subsection{Partición por CSV/día}
La partición por CSV/día entrena sobre ficheros seleccionados de CICIDS2017 y evalúa sobre ficheros diferentes. Esta partición comprueba si el modelo se transfiere entre días de captura y contextos de ataque. En la implementación mantenida, la partición por día por defecto entrena sobre los patrones del lunes, martes y miércoles, y prueba sobre los patrones del jueves y viernes. Se trata de una configuración más estricta que la partición aleatoria porque ficheros enteros quedan fuera del entrenamiento.

La partición por día también se alinea con la estructura conocida de CICIDS2017: diferentes días contienen diferentes cargas de trabajo benignas y escenarios de ataque. Por ello es más adecuada para revelar si el modelo ha aprendido un comportamiento de flujo general o artefactos específicos de un día.

\subsection{Partición leave-one-exact-CSV-out}
El protocolo interno más estricto es la validación leave-one-exact-CSV-out. Cada fold retiene un CSV oficial de CICIDS2017 como datos de prueba y entrena con los CSV restantes. El script de validación implementado agrega métricas por fold, incluyendo exactitud, exactitud balanceada, precisión de ataque, recall de ataque, F1 de ataque, precisión de benigno, recall de benigno, tasa de falsos positivos, tasa de falsos negativos, tasa de bloqueo, retorno total y costes temporales.

Este protocolo es metodológicamente importante porque trata cada fichero de captura como una unidad de tipo dominio. Es más sólido que una partición por día basada en subcadenas porque utiliza los nombres de fichero CSV oficiales exactos y evita solapamientos accidentales. En el momento de redacción de este capítulo, la implementación existe, mientras que las métricas agregadas finales solo deben reportarse cuando se disponga de un artefacto de ejecución comprometido.

\subsection{Estandarización ajustada sobre entrenamiento}
El entrenamiento de redes neuronales es sensible a la escala de las características. El pipeline utiliza por tanto normalización mediante \texttt{StandardScaler}, ajustando el escalador únicamente sobre la partición de entrenamiento y aplicando la transformación ajustada a los datos de prueba. El mismo principio se aplica a los artefactos de inferencia persistidos: el escalador guardado durante el entrenamiento se reutiliza en la Fase 2, en lugar de recalcularse a partir de los flujos de laboratorio.

Para la inferencia robusta en la Fase 2, el pipeline de entrenamiento también persiste los límites de percentiles sobre las características canónicas en bruto. El script de inferencia de la Fase 2 puede aplicar recorte por percentiles sobre las características en bruto y recorte por z-score después del escalado, y luego exportar diagnósticos. Estos controles no se presentan como una solución al cambio de dominio; son salvaguardas que hacen visibles y acotados los valores extremos fuera de distribución durante la inferencia offline.

\section{Formulación del aprendizaje por refuerzo}
El defensor se implementa como un entorno compatible con Gymnasium \parencite{GymnasiumDocs}. Esto permite que el código estándar de RL basado en valores interactúe con el conjunto de datos de flujos a través de la interfaz habitual \texttt{reset()} y \texttt{step()}.

\subsection{Espacio de observación}
El espacio de observación es un espacio vectorial continuo de dimensión 152. Cada observación representa un vector de flujo canónico escalado junto con su máscara de presencia/ausencia. Durante el entrenamiento, las observaciones se extraen de la partición de entrenamiento. Durante la evaluación determinista, el entorno itera a través de la partición de prueba sin barajado, lo que permite comparar las predicciones directamente con las etiquetas verdaderas.

\subsection{Espacio de acciones}
El espacio de acciones es discreto con dos acciones:
\begin{itemize}
    \item \texttt{Acción 0 (PERMIT)}: clasificar el flujo como benigno y permitirlo en el modelo de decisión experimental;
    \item \texttt{Acción 1 (BLOCK)}: clasificar el flujo como malicioso y bloquearlo o generar una alerta en el modelo de decisión experimental.
\end{itemize}

El valor de la acción está alineado intencionalmente con el valor de la etiqueta binaria: benigno es \texttt{0}, ataque es \texttt{1}. Esto simplifica la evaluación directa porque una acción predicha puede compararse con la etiqueta objetivo sin necesidad de una capa de decodificación adicional.

\subsection{Función de recompensa}
La función de recompensa codifica costes de seguridad asimétricos \parencite{Lee2002CostSensitiveIDS,Cardenas2006IDSFramework}. Un falso negativo significa que se permite un ataque. Un falso positivo significa que se bloquea tráfico benigno. En muchos escenarios defensivos, los ataques no detectados son más graves que los bloqueos innecesarios, aunque la relación exacta depende del contexto operativo.

La implementación actual utiliza la siguiente configuración de recompensa por defecto:
\[
\text{TP}=1.5,\quad \text{FP}=-2.0,\quad \text{FN}=-5.0,\quad \text{TN/omisión}=0.0.
\]
La Tabla~\ref{tab:reward-matrix} organiza estos valores como una matriz de decisión para los cuatro resultados posibles.

\begin{table}[htbp]
\centering
\begin{tabular}{lcc}
\toprule
Decisión del agente & Etiqueta real: Ataque & Etiqueta real: Benigno \\
\midrule
\texttt{BLOCK} (Acción 1) & $+1.5$ (Verdadero Positivo) & $-2.0$ (Falso Positivo) \\
\texttt{PERMIT} (Acción 0) & $-5.0$ (Falso Negativo) & $\phantom{-}0.0$ (Verdadero Negativo) \\
\bottomrule
\end{tabular}
\caption{Matriz de recompensa por defecto para el entorno del defensor binario. Los valores son supuestos experimentales registrados con cada ejecución.}
\label{tab:reward-matrix}
\end{table}

La penalización por falso negativo domina la señal de recompensa. Un agente que emite \texttt{PERMIT} en cada observación acumula un retorno muy negativo en tráfico con alta proporción de ataques, mientras que un agente que emite \texttt{BLOCK} en cada observación acumula penalizaciones por falsos positivos sin obtener crédito por detección de ataques. La estructura de recompensa está por tanto diseñada para alejar al agente de políticas constantes degeneradas y conducirlo hacia un comportamiento discriminativo. La recompensa cero para el verdadero negativo es deliberada: recompensar explícitamente cada paso benigno correcto sesgaría al agente hacia maximizar la tasa de permiso en lugar de la calidad de detección. Estos valores son supuestos experimentales, no costes empresariales medidos. Se registran con cada ejecución para que las ejecuciones históricas puedan distinguirse de los valores por defecto actuales.

\subsection{Estructura del episodio}
Un episodio es una secuencia de decisiones sobre flujos. Al reiniciar, el entorno inicializa un índice sobre la partición seleccionada. Cuando el barajado está activado, el orden de entrenamiento se aleatoriza. En cada paso, el agente recibe una observación, selecciona \texttt{PERMIT} o \texttt{BLOCK}, recibe una recompensa calculada a partir de la etiqueta real y la acción, y avanza al flujo siguiente. Un episodio termina cuando se agota el conjunto de datos o se alcanza un número máximo de pasos configurado.

Este diseño proporciona a los algoritmos estándar de RL un flujo de transiciones manteniendo el entrenamiento offline y seguro. Ninguna acción afecta al conjunto de datos fuente, al flujo siguiente ni al entorno de red.

\subsection{Estado del entorno y protocolo de transición}
En cada llamada a \texttt{reset()}, el entorno selecciona la partición asignada al modo actual (entrenamiento o evaluación), baraja opcionalmente el índice si el modo de barajado está activo, reinicia el contador de pasos interno y devuelve la primera observación junto con un diccionario de metadatos. En cada llamada a \texttt{step()}, el entorno registra la acción del agente, recupera la etiqueta de verdad del índice de flujo actual, calcula la recompensa escalar a partir de la matriz de recompensa, avanza el índice y devuelve la siguiente observación, la recompensa, un indicador de terminación, un indicador de truncación y el diccionario de metadatos.

La condición de terminación es el agotamiento del conjunto de datos o un número máximo de pasos configurable. En el modo de entrenamiento, el orden barajado garantiza que las transiciones del minibatch muestreen una distribución amplia de tipos de flujo durante un único episodio. En el modo de evaluación, el entorno es determinista: la misma partición recorrida en el mismo orden produce la misma secuencia de observaciones, recompensas y etiquetas. Este determinismo es esencial para la evaluación reproducible porque implica que las diferencias en métricas entre dos modelos guardados reflejan únicamente el comportamiento del modelo y no el ruido de muestreo en el bucle de evaluación.

La separación entre los modos de entrenamiento y evaluación se aplica a nivel del entorno y no se deja al llamante. Pasar el indicador de modo correcto en el momento de la construcción establece tanto el conjunto de índices como el comportamiento de barajado, lo que evita un error frecuente en el que la evaluación hereda accidentalmente el orden barajado de un entorno de entrenamiento.

\subsection{Limitaciones de la formulación conjunto-de-datos-como-entorno}
Dado que el entorno es un conjunto de datos estático, las acciones del agente no afectan a los estados futuros. Esto no es un problema completo de defensa cibernética autónoma en el que bloquear un flujo podría cambiar el comportamiento posterior del atacante, el estado del host, el enrutamiento o el contexto de alertas. La formulación se asemeja más a un problema de decisión contextual sensible al coste que a un proceso de decisión de Markov rico \parencite{SuttonBarto2018RL,Levine2020OfflineRL}.

Esta limitación es explícita en la metodología. El RL se utiliza aquí para estudiar el aprendizaje de decisiones basadas en valor y sensibles al coste sobre registros de flujo, y para establecer un pipeline offline seguro. Las afirmaciones sobre respuesta adversarial en múltiples pasos, adaptación en línea o defensa cibernética activa quedan fuera del alcance implementado.

\section{Agente QRDQN}
El algoritmo de RL principal es Quantile Regression Deep Q-Network (QRDQN) \parencite{Dabney2018QRDQN}. QRDQN pertenece a la familia del RL distribucional, que modela una distribución sobre los retornos en lugar de únicamente un retorno esperado \parencite{Bellemare2017DistributionalRL}.

\subsection{Papel algorítmico}
El DQN estándar aproxima los valores-acción con una red neuronal y selecciona acciones según el retorno esperado estimado \parencite{Mnih2015DQN}. QRDQN extiende esta idea aprendiendo los cuantiles de la distribución del retorno. Esto es relevante para la tesis porque la estructura de recompensa es asimétrica: los falsos negativos conllevan una penalización mucho mayor que la recompensa que generan los permisos benignos correctos. Una visión distribucional puede representar más información sobre la variabilidad del retorno que una única estimación de la media.

La metodología no asume que QRDQN sea inherentemente superior a los clasificadores supervisados para datos de flujo tabulares. En cambio, evalúa QRDQN como el representante principal de RL bajo el mismo preprocesamiento, las mismas particiones y las mismas métricas utilizadas por las líneas base.

\subsection{Función de pérdida de regresión cuantílica}
QRDQN representa la distribución del retorno de acción mediante $N$ estimaciones de cuantiles aprendidas $\{\theta_i(s, a)\}_{i=1}^{N}$ en puntos medios de cuantil fijos $\hat{\tau}_i = \frac{2i-1}{2N}$ \parencite{Dabney2018QRDQN}. Durante la selección de acción, estos cuantiles se promedian para producir una estimación del retorno esperado, recuperando la regla codiciosa estándar del DQN bajo expectativa.

Para una única transición $(s, a, r, s')$, sea $a^* = \arg\max_a \sum_j \theta_j^-(s', a)$ la acción codiciosa bajo la red objetivo, y sea
\[
    \delta_{ij} = r + \gamma\,\theta_j^-(s', a^*) - \theta_i(s, a)
\]
el error de diferencia temporal entre el cuantil fuente $i$ y el cuantil objetivo $j$ calculado mediante bootstrapping. QRDQN minimiza la función de pérdida de regresión cuantílica de Huber sumada sobre todos los $N^2$ pares de cuantiles:
\[
    \mathcal{L} = \frac{1}{N}\sum_{i=1}^{N}\sum_{j=1}^{N}
        \bigl|\hat{\tau}_i - \mathbf{1}[\delta_{ij} < 0]\bigr|\cdot L_\kappa(\delta_{ij}),
\]
donde $L_\kappa$ es la pérdida de Huber elemento a elemento con umbral $\kappa$:
\[
    L_\kappa(\delta) = \begin{cases} \tfrac{1}{2}\delta^2 & \text{si } |\delta| \leq \kappa, \\ \kappa\bigl(|\delta| - \tfrac{1}{2}\kappa\bigr) & \text{en caso contrario.}\end{cases}
\]
A diferencia de C51, que representa los retornos sobre un soporte categórico fijo, QRDQN no impone ninguna restricción sobre el soporte de la distribución del retorno y se adapta automáticamente a la escala de las recompensas \parencite{Bellemare2017DistributionalRL}. En el entorno de recompensa asimétrica de esta tesis, la distribución del retorno puede desarrollar colas negativas pronunciadas cuando los eventos de falso negativo son frecuentes, lo que hace que una visión distribucional sea más informativa que una estimación escalar de la media para comprender el comportamiento del agente en torno a la frontera de decisión PERMIT/BLOCK.

\subsection{Factor de descuento $\gamma = 0$: formulación de un solo paso}
En todas las ejecuciones de esta tesis, incluida la ejecución oficial, el factor de descuento se fija en $\gamma = 0$ (valor registrado en el artefacto de configuración de la ejecución). La elección es deliberada y coherente con la formulación conjunto-de-datos-como-entorno: dado que el entorno es un conjunto de datos estático y las acciones del agente no afectan a los estados futuros, no existe dependencia temporal entre flujos que un factor de descuento positivo pudiera explotar.

Con $\gamma = 0$, el error de diferencia temporal definido más arriba se reduce a
\[
    \delta_{ij} = r - \theta_i(s, a),
\]
es decir, el término de bootstrapping $\gamma\,\theta_j^-(s', a^*)$ se anula y la red objetivo deja de contribuir. Cada cuantil aprendido $\theta_i(s,a)$ regresa por tanto hacia la recompensa inmediata de la decisión, y QRDQN actúa como un estimador distribucional de un solo paso: formalmente, un \emph{bandit contextual sensible al coste}, y no un proceso de decisión de Markov secuencial.

Esta reducción no vacía de contenido la elección de QRDQN. La cabeza distribucional sigue modelando la distribución completa de la recompensa asimétrica en torno a la frontera \texttt{PERMIT}/\texttt{BLOCK} —y no únicamente su media—, lo que aporta información sobre la incertidumbre del valor bajo coste asimétrico. La diferencia frente a un clasificador supervisado no reside en la presencia de dinámica secuencial, inexistente aquí por diseño, sino en que el aprendizaje se realiza mediante valores sobre una matriz de coste explícita (con la penalización de falso negativo muy superior a la de falso positivo) en lugar de mediante entropía cruzada sobre etiquetas. La tesis no reivindica autonomía secuencial ni respuesta adversarial en múltiples pasos: tal y como se recoge en las limitaciones de la formulación, esas capacidades quedan fuera del alcance implementado.

\subsection{Estrategia de exploración}
Durante el entrenamiento, el agente emplea una política $\varepsilon$-greedy. Con probabilidad $\varepsilon$ selecciona una acción uniformemente aleatoria de $\{\texttt{PERMIT}, \texttt{BLOCK}\}$; con probabilidad $1-\varepsilon$ selecciona la acción cuya estimación media de cuantiles es más alta \parencite{Mnih2015DQN}. $\varepsilon$ decae linealmente desde un valor inicial cercano a uno hasta un valor final pequeño a lo largo de una fracción fija del total de pasos de entrenamiento, y luego permanece constante. En el entorno de acción binaria, el presupuesto de exploración es relativamente barato porque solo existen dos alternativas. La fracción de exploración y el valor final de $\varepsilon$ se almacenan en el artefacto de configuración de la ejecución para que las ejecuciones con diferentes programas no se mezclen o comparen en silencio sin tener en cuenta sus diferencias de exploración.

\subsection{Configuración del entrenamiento}
La implementación utiliza el módulo QRDQN de \texttt{sb3-contrib} \parencite{SB3ContribQRDQNDocs}. El modelo emplea una política de perceptrón multicapa sobre las observaciones de 152 dimensiones. El entrenamiento utiliza un buffer de repetición, actualizaciones por minibatch, una red objetivo y predicción determinista durante la evaluación. El script de entrenamiento mantenido admite preajustes rápidos y completos, modos de partición aleatoria y por día, semillas explícitas, límites de filas configurables y pasos temporales configurables.

Los valores por defecto del preajuste completo actual utilizan lotes más grandes y más pasos de gradiente que el preajuste rápido. Esta separación permite al proyecto ejecutar comprobaciones rápidas de humo sin confundirlas con las ejecuciones finales del benchmark. Cualquier resultado reportado debe incluir el identificador exacto de la ejecución y la configuración.

\subsection{Persistencia del modelo y el preprocesamiento}
El entrenamiento produce un directorio de ejecución y un artefacto de modelo. El directorio de ejecución almacena la configuración, las métricas, el escalador ajustado y los límites de percentiles de entrenamiento. Persistir los artefactos de preprocesamiento es esencial porque la inferencia en la Fase 2 debe utilizar el mismo estado de normalización que el entrenamiento. Recalcular un escalador sobre el tráfico de laboratorio filtraría información de la distribución de inferencia hacia la interfaz del modelo y haría la ejecución irreproducible.

\section{Línea base supervisada}
Un método de RL para NIDS a nivel de flujo debe compararse con líneas base supervisadas sólidas. Los datos de flujo tabulares son un entorno donde los modelos basados en árboles, especialmente los Random Forests y los ensamblados con gradient boosting, suelen ser competitivos \parencite{LiuLang2019IDSSurvey,Apruzzese2022CrossEvaluationNIDS}.

\subsection{Línea base con Random Forest}
La línea base supervisada principal es un clasificador Random Forest. La línea base es metodológicamente útil porque obtiene buenos resultados en datos tabulares, gestiona las interacciones no lineales entre características y requiere menos complejidad de entrenamiento que el RL profundo. También proporciona una comprobación de cordura frente a la sobreestimación del valor de la formulación de RL en una tarea que puede comportarse como una clasificación sensible al coste.

\subsection{Comparación bajo el mismo protocolo}
La línea base debe consumir la misma representación canónica de observaciones, utilizar los mismos protocolos de partición y reportar las mismas métricas que QRDQN. Esta regla de mismo protocolo evita comparaciones sesgadas. Si la línea base utiliza un preprocesamiento diferente, controles de fuga distintos o particiones más sencillas, la comparación no aísla la contribución algorítmica de QRDQN.

La metodología trata por tanto la comparación con la línea base como un requisito del protocolo y no como un apéndice opcional. También evita afirmar la superioridad de QRDQN hasta que existan artefactos de línea base con la misma partición.

\subsection{Desbalanceo de clases y ponderación sensible al coste}
CICIDS2017 presenta desbalanceo de clases: los flujos benignos superan sustancialmente en número a los flujos de ataque, y dentro de los flujos de ataque las familias individuales varían en prevalencia. Un Random Forest entrenado con pesos de clase uniformes por defecto tenderá por tanto a sesgar sus predicciones hacia la clase benigna mayoritaria. Para que la comparación sea significativa bajo los supuestos de coste asimétrico de la tesis, la línea base de Random Forest utiliza ponderación de clases balanceada, donde cada clase recibe un peso inversamente proporcional a su frecuencia en la partición de entrenamiento. Esto garantiza que el modelo supervisado no sea penalizado en la misma dirección que la función de recompensa de RL por razones diferentes.

La configuración de los pesos de clase se registra en el artefacto de configuración de la línea base junto al modo de partición, el esquema de características y la semilla de evaluación, de modo que el supuesto de ponderación es trazable. Si la tesis encuentra que el Random Forest con ponderación balanceada alcanza un rendimiento en falsos negativos similar o superior al de QRDQN bajo la recompensa asimétrica, esa comparación es informativa precisamente porque ambos modelos han tenido acceso a mecanismos de sensibilidad al coste comparables. Comparar un QRDQN ajustado al coste frente a un Random Forest sin corrección no aportaría esta evidencia.

\section{Protocolo de escala de entrenamiento y eficiencia de datos}
La metodología incluye un componente de escala de entrenamiento motivado por una pregunta práctica: ¿cuántos datos de flujo etiquetados necesita cada enfoque antes de alcanzar un rendimiento estable? Esto es relevante porque el tráfico de ataque etiquetado es costoso de obtener, la abundancia de benchmarks no refleja la disponibilidad en el despliegue, y los métodos de RL profundo pueden ser sustancialmente más exigentes en cuanto a muestras que los clasificadores supervisados basados en árboles \parencite{SuttonBarto2018RL,Ring2019DatasetSurvey}.

\subsection{Protocolo de curva de aprendizaje}
El protocolo de curva de aprendizaje evalúa tanto el agente QRDQN como la línea base de Random Forest con presupuestos de entrenamiento progresivamente mayores, todos extraídos de la misma partición estratificada. El conjunto de evaluación se mantiene constante en todos los presupuestos como la partición de prueba completa, de modo que los cambios en las métricas reflejan únicamente el tamaño de los datos de entrenamiento y no las condiciones de evaluación. Las métricas se registran en cada nivel de presupuesto utilizando el mismo código de evaluación que las ejecuciones principales del benchmark.

Cada presupuesto utiliza un escalador ajustado de nuevo calculado únicamente a partir de las filas incluidas en ese subconjunto de entrenamiento. Los presupuestos más grandes se benefician por tanto de una normalización más representativa sin contaminar la evaluación fija de prueba. Esto significa que los artefactos de preprocesamiento producidos en cada nivel de presupuesto son distintos e independientemente reproducibles.

\subsection{Niveles de presupuesto y comparación}
Los niveles de presupuesto candidatos incluyen 100\,k, 250\,k, 500\,k, 1\,M y 2\,M registros de flujo. No todos los niveles son alcanzables bajo todos los modos de partición; cuando una partición no produce conjuntos de entrenamiento del tamaño objetivo, se registra el tamaño disponible más cercano. El presupuesto final es la partición de entrenamiento completa bajo la partición seleccionada.

Comparar QRDQN y Random Forest en cada nivel de presupuesto revela si la formulación de RL tiene un requisito de muestras o una trayectoria de convergencia diferente. Si el Random Forest alcanza una meseta de rendimiento estable con sustancialmente menos filas, esto constituye evidencia interna de que la tarea del benchmark no requiere la complejidad de muestras adicional del RL profundo sobre el esquema de características actual. Por el contrario, si el comportamiento de RL sensible al coste mejora desproporcionadamente a medida que crece el presupuesto, los resultados de escala de entrenamiento pueden identificar a qué escala la estructura de recompensa distribucional comienza a aportar una ventaja medible. Todos los resultados de escala de entrenamiento se almacenan bajo el mismo esquema de artefactos de ejecución que los experimentos principales, etiquetados por nivel de presupuesto y semilla, de modo que permanecen recuperables de forma independiente y comparables en experimentos futuros.

\section{Salidas de evaluación y reproducibilidad}
La reproducibilidad es una restricción de diseño fundamental, especialmente porque la investigación en seguridad con ML suele adolecer de código faltante, detalles de procesamiento de datos incompletos y registros de evaluación incompletos \parencite{Olszewski2023ReproSecurityML}.

\subsection{Métricas}
Las métricas binarias principales son exactitud, precisión, recall y F1 tanto para la clase benigna como para la de ataque. El análisis también rastrea los conteos de la matriz de confusión: verdaderos positivos, falsos positivos, falsos negativos y verdaderos negativos. Estos conteos son necesarios porque dos modelos con exactitud agregada similar pueden tener implicaciones de seguridad muy diferentes.

La tasa de falsos positivos y la tasa de falsos negativos son especialmente importantes. La tasa de falsos positivos mide la proporción de tráfico benigno bloqueado incorrectamente, lo que se relaciona con la disponibilidad y la carga de trabajo del analista. La tasa de falsos negativos mide la proporción de ataques permitidos incorrectamente, lo que se relaciona directamente con el riesgo de intrusiones no detectadas. El análisis de precisión y recall es particularmente relevante bajo desbalanceo de clases, donde la exactitud sola puede ser engañosa \parencite{Fawcett2006ROC,Saito2015PRROC}.

\subsection{Escalera de validación}
La escalera de evaluación está ordenada de la evidencia más débil a la más sólida:
\begin{enumerate}
    \item partición aleatoria estratificada, utilizada como comprobación de la capacidad de aprendizaje;
    \item comprobación de predicción directa, que verifica las salidas del modelo frente a las etiquetas de prueba sin depender de los metadatos del entorno;
    \item comprobación de etiquetas barajadas, que ayuda a detectar fugas de información confirmando que el rendimiento se degrada cuando las etiquetas de entrenamiento se aleatorizan;
    \item partición por CSV/día, que retiene grupos de captura completos;
    \item validación leave-one-exact-CSV-out, que retiene cada CSV oficial por turnos;
    \item inferencia sobre tráfico de laboratorio en la Fase 2, que prueba el pipeline entrenado sobre tráfico capturado de forma independiente.
\end{enumerate}

En la implementación, los peldaños segundo a cuarto se corresponden con las comprobaciones \emph{Check A} (predicción directa), \emph{Check B} (etiquetas barajadas) y \emph{Check C} (partición por CSV/día) del script \texttt{src/validate\_checks.py}; el quinto peldaño se implementa en \texttt{src/validate\_leave\_one\_csv\_out.py}.

Cada peldaño apunta a un modo de fallo diferente. La escalera no hace que ningún resultado individual sea definitivo, pero hace la evaluación más informativa que una única puntuación aleatoria de entrenamiento-prueba.

La partición aleatoria estratificada comprueba si el modelo y la implementación pueden resolver el benchmark en absoluto. Un modelo que falla aquí tiene un problema de entrenamiento o de arquitectura, no un problema de generalización. La comprobación de etiquetas barajadas es un detector de fugas: un modelo que alcanza una exactitud muy por encima de la de la clase mayoritaria en una ejecución con etiquetas barajadas ha aprendido a partir de artefactos correlacionados con las etiquetas en las características y no a partir de estadísticas de flujo reales \parencite{Arp2020DosDontsMLSecurity}. La partición por CSV/día comprueba la generalización temporal y por escenario reteniendo grupos de captura completos; la degradación del rendimiento respecto a la partición aleatoria revela si el modelo depende del contexto específico del día. El protocolo leave-one-exact-CSV-out va más lejos: cada fichero oficial de CICIDS2017 se retiene por turnos, y la distribución de métricas agregada muestra si algún fichero domina el benchmark o si el rendimiento es estable entre escenarios de captura. La inferencia en la Fase 2 es el peldaño final y más estricto, que comprueba si todo el pipeline puede operar sobre tráfico producido fuera del entorno del benchmark.

\subsection{Artefactos de ejecución}
Toda ejecución de entrenamiento o evaluación significativa debe persistir un \texttt{RUN\_ID} único, un fichero de configuración, métricas y cualquier estado de preprocesamiento necesario para su reutilización exacta. Las ejecuciones de entrenamiento persisten los pesos del modelo y los artefactos del escalador. Las ejecuciones de validación persisten los resultados de validación. Las ejecuciones de inferencia de la Fase 2 persisten las predicciones, las métricas, la configuración y los diagnósticos opcionales.

Esta disciplina de artefactos evita afirmaciones ambiguas. Un resultado solo se trata como medido cuando puede vincularse a una carpeta de ejecución y una configuración específicas. Las ejecuciones históricas no se fusionan en silencio con los valores por defecto actuales; los valores de recompensa, la lógica de partición, el preprocesamiento, los límites de filas y las semillas deben leerse del artefacto.

\section{Pipeline de inferencia offline en la Fase 2}
El punto de entrada principal de la Fase 2 mantenido es el script de inferencia offline robusta. Acepta un CSV de flujos extraído, un modelo entrenado, un escalador persistido, límites de percentiles opcionales y recorte por z-score opcional. A continuación mapea las columnas de entrada al esquema canónico, aplica la misma representación de características y máscara, escala con el escalador guardado, calcula diagnósticos, predice acciones en lotes y escribe las salidas en un nuevo directorio de ejecución de la Fase 2.

\subsection{Mapeo canónico para flujos de laboratorio}
Los extractores de flujos de laboratorio pueden utilizar nombres de columna diferentes a los de CICIDS2017. El script de la Fase 2 define por tanto un mapeo explícito de los nombres del extractor a los nombres canónicos. Las columnas de metadatos como la IP de origen, la IP de destino, los puertos, el protocolo y la marca temporal pueden preservarse en los ficheros de salida para su inspección, pero no forman parte del vector de características del modelo. Esto mantiene separadas la interpretabilidad y la entrada al modelo.

\subsection{Diagnósticos de cambio de distribución}
La inferencia de la Fase 2 calcula diagnósticos sobre las características canónicas escaladas. Los z-scores absolutos elevados indican que los valores de flujo de laboratorio están lejos de la distribución de entrenamiento de CICIDS2017. El script también puede exportar las características con las mayores desviaciones. Estos diagnósticos no son etiquetas y no prueban la corrección, pero ayudan a explicar los cambios en la tasa de bloqueo e identificar discrepancias entre el tráfico del benchmark y el del laboratorio.

Más allá de los z-scores por característica, la salida de diagnóstico registra la proporción de características canónicas que fueron imputadas a partir de la máscara de presencia/ausencia, la media y la varianza de los valores de características escaladas relativas a las estadísticas de entrenamiento, y cualquier característica que caiga consistentemente fuera de los límites de percentiles observados durante el entrenamiento con CICIDS2017. En conjunto, estos indicadores proporcionan un resumen estructurado de la distancia distribucional entre el tráfico de laboratorio y el benchmark, permitiendo una investigación dirigida en lugar de requerir la inspección manual de los registros brutos de predicción.

\subsection{Métricas de verdad cuando están disponibles}
Si un CSV de flujos de la Fase 2 incluye columnas de verdad como \texttt{truth\_y} o \texttt{truth\_label}, el script puede calcular las mismas métricas binarias utilizadas en la evaluación interna. Si no se dispone de la verdad fundamental, el script reporta estadísticas de predicción como la tasa de bloqueo, la tasa de permiso y los diagnósticos de z-score. Esta distinción es importante: una ejecución de laboratorio no etiquetada, solo benigna o mixta, no puede respaldar las mismas afirmaciones que un conjunto de validación externo etiquetado.

\section{Marco de implementación}
El pipeline experimental está implementado en Python y se apoya en un conjunto de bibliotecas de código abierto bien establecidas. El componente de RL profundo principal utiliza la extensión \texttt{sb3-contrib} de Stable Baselines3, que proporciona una implementación contrastada de QRDQN con arquitectura de red configurable, buffer de repetición, programa de exploración y actualizaciones de la red objetivo \parencite{SB3ContribQRDQNDocs}. El envoltorio del entorno sigue la API de Gymnasium \parencite{GymnasiumDocs}, haciendo que el conjunto de datos de flujos sea compatible con cualquier bucle de entrenamiento que cumpla con Gymnasium sin cambios en la implementación del algoritmo.

\subsection{Pila de procesamiento de datos y aprendizaje supervisado}
La carga y el preprocesamiento de datos se apoyan en \texttt{pandas} para la ingesta de CSV, la estandarización de columnas y la lectura fragmentada de ficheros grandes. Las matrices de características se convierten a arrays de \texttt{numpy} y se castean a \texttt{float32} antes de entrar en el entorno de Gymnasium o en la línea base supervisada. La estandarización utiliza \texttt{StandardScaler} de \texttt{scikit-learn}, elegido por su compatibilidad con el protocolo de serialización de \texttt{sklearn}, que permite persistir y recargar los escaladores ajustados sin reimplementación. La línea base de Random Forest también utiliza \texttt{scikit-learn}, consumiendo los mismos arrays de entrada que entran en el entorno de RL. Esta ruta de datos compartida garantiza que cualquier diferencia de rendimiento entre los dos modelos refleje el comportamiento algorítmico y no una discrepancia en el manejo de los datos.

\subsection{Gestión de artefactos}
Cada ejecución de entrenamiento produce un \texttt{RUN\_ID} único compuesto por una marca temporal y un sufijo aleatorio corto. El directorio de ejecución reúne la política QRDQN serializada, el escalador ajustado, los límites de percentiles, un JSON de configuración y un JSON de métricas. Las ejecuciones de validación y de inferencia de la Fase 2 producen subdirectorios independientes con sus propios registros de configuración y métricas. Este esquema de directorios plano evita dependencias de bases de datos al tiempo que hace las ejecuciones recuperables de forma independiente y comparables mediante inspección del sistema de ficheros.

\subsection{Controles de reproducibilidad}
Las semillas se fijan para los generadores de números aleatorios de \texttt{numpy}, \texttt{random} y \texttt{torch} al comienzo de cada ejecución. Los reinicios del entorno de Gymnasium también reciben una semilla explícita. Cuando se utilizan múltiples semillas para la estimación de la varianza, cada semilla genera un directorio de ejecución independiente de modo que los resultados por semilla se preservan individualmente en lugar de promediarse en silencio. Estos controles reducen la varianza interna a la ejecución lo suficiente como para distinguir diferencias de rendimiento significativas del ruido, aunque no pueden eliminar todo el indeterminismo derivado del ordenamiento de los kernels de la GPU.

\section{Limitaciones metodológicas}
La metodología realiza compromisos deliberados que definen la interpretación correcta de los resultados de la tesis. Esta sección organiza esos compromisos en cuatro categorías.

\subsection{Entorno estático y toma de decisiones secuencial}
La limitación más fundamental es que el entorno de RL es un conjunto de datos estático. Las acciones del agente no alteran las observaciones de flujo futuras, la topología de la red, el comportamiento del atacante ni el estado del sistema. La formulación se asemeja por tanto más a un problema de decisión contextual sensible al coste que a un proceso de decisión de Markov completo \parencite{SuttonBarto2018RL,Levine2020OfflineRL}. Las dependencias secuenciales, las cadenas de ataque de múltiples etapas, la adaptación del atacante y las contramedidas activas del defensor quedan todas fuera del alcance implementado. Las afirmaciones sobre autonomía secuencial, defensa cibernética adaptativa o aprendizaje por refuerzo en línea no están respaldadas por el pipeline actual.

\subsection{Mapeo de etiquetas binarias y detalle por familia de ataque}
El mapeo de etiquetas binarias colapsa todos los tipos de ataque en una única clase \texttt{BLOCK}. La información sobre la familia de ataque se descarta en la ruta de aprendizaje principal de RL. Este diseño está justificado por el modelo de decisión binario estudiado aquí, pero significa que el agente no puede especializar su calidad de detección para familias de ataque individuales. El análisis de errores por familia requiere etiquetas de ataque preservadas y código de evaluación independiente, y no puede inferirse a partir de las métricas de exactitud binaria reportadas por el pipeline principal.

\subsection{Fragilidad del benchmark y generalización}
CICIDS2017 sigue siendo un benchmark público controlado con problemas documentados: artefactos de extracción, inconsistencias en las etiquetas, flujos duplicados y sensibilidad a la partición \parencite{Engelen2021CICIDSIssues,Lanvin2023Errors,Liu2022ErrorPrevalenceNIDS}. Incluso con preprocesamiento anti-fuga y protocolos de partición más estrictos, el rendimiento sobre CICIDS2017 no equivale al rendimiento sobre una red de producción. La inferencia de laboratorio en la Fase 2 comprueba el cambio de dominio, pero es inferencia offline sobre una captura privada y no una validación sobre un flujo de producción etiquetado de forma independiente. Una evidencia de generalización más sólida requeriría conjuntos de datos etiquetados externos, múltiples entornos de captura y, en última instancia, un análisis de seguridad para cualquier despliegue activo.

\subsection{Algoritmo único y alcance de la reproducibilidad}
El algoritmo de RL principal es una única configuración de QRDQN. Los resultados de una sola ejecución no pueden cuantificar la varianza entre ejecuciones a menos que se evalúen múltiples semillas y sus resultados se agreguen \parencite{SuttonBarto2018RL}. La línea base supervisada se limita a Random Forest en la comparación principal; otros modelos tabulares se discuten en la revisión de la literatura pero no forman parte del protocolo evaluado. Estas restricciones de alcance son prácticas: completar el pipeline central y la escalera de validación es metodológicamente más informativo que implementar parcialmente muchas comparaciones. Afirmaciones más sólidas en el futuro requerirían artefactos completos de leave-one-exact-CSV-out, métricas de línea base supervisada con el mismo protocolo, múltiples semillas, tráfico de laboratorio etiquetado más rico y, en última instancia, un análisis de seguridad independiente para cualquier despliegue activo.

\chapter{Resultados}

Este capítulo separa el estado de la evidencia de su interpretación. En la fecha de este borrador, la campaña final definida en \path{experiments/final_experiment_campaign.json} todavía no se ha ejecutado. En consecuencia, primero se fija la estructura en la que se incorporarán sus artefactos y, después, se conserva una síntesis de resultados históricos útiles para supervisión. Estos últimos no constituyen la evaluación final del sistema y serán sustituidos o relegados cuando exista una campaña completa y validada.

\section{Condiciones generales y trazabilidad}
\label{sec:resultados-condiciones}

Una cifra se considera evidencia final únicamente si puede trazarse desde una ejecución de la campaña aprobada hasta su \texttt{config.json}, \texttt{metrics.json} o fichero de validación, \texttt{environment.json}, resumen, manifiesto y sumas de comprobación. La Figura~\ref{fig:trazabilidad-evidencia} resume esa cadena. Esta regla evita mezclar tres objetos distintos: el perfil experimental congelado, los artefactos históricos ya existentes y los artefactos finales aún pendientes.

\begin{figure}[htbp]
\centering
\resizebox{0.98\textwidth}{!}{% Figura conceptual: cadena de evidencia, sin resultados medidos.
% Fuente: esquema de artefactos 3.0 y agregador de campaña.
\begin{tikzpicture}[node distance=3mm]
  \node[cajaDato, text width=2.1cm] (spec) {Diseño congelado\\+ perfil hash};
  \node[cajaProceso, text width=2.1cm, right=of spec] (run) {Ejecución física\\+ RUN\_ID};
  \node[cajaProceso, text width=2.1cm, right=of run] (art) {Artefactos 3.0\\+ SHA-256};
  \node[cajaModelo, text width=2.1cm, right=of art] (agg) {Agregado\\validado};
  \node[cajaResaltada, text width=2.1cm, right=of agg] (thesis) {Tabla / figura /\\afirmación};

  \draw[flechaAzul] (spec) -- (run);
  \draw[flechaAzul] (run) -- (art);
  \draw[flechaAzul] (art) -- (agg);
  \draw[flechaAzul] (agg) -- (thesis);

  \node[cajaPendiente, text width=8.6cm, below=7mm of art] (gate)
    {\textbf{Puerta de aceptación:} campaña completa, hashes coherentes, dependencias correctas, métricas definidas o nulas y ausencia de sustitutos históricos};
  \draw[flechaFina] (art) -- (gate);
  \draw[flechaFina] (gate.east) -| (agg.south);
\end{tikzpicture}
}
\caption[Trazabilidad desde el diseño hasta la tesis]{Cadena de trazabilidad exigida para incorporar una cifra final: especificación congelada, ejecución identificada, artefactos sellados, agregación validada y elemento de tesis. Una ausencia o inconsistencia detiene la incorporación. Fuente: elaboración propia a partir del esquema de artefactos de campaña.}
\label{fig:trazabilidad-evidencia}
\end{figure}

La Tabla~\ref{tab:estructura-resultados-finales} fija el orden definitivo. No se crean tablas vacías para simular resultados: cada bloque se incorporará cuando el agregador haya comprobado que están presentes todas las ejecuciones físicas y alias previstos.

\begin{table}[htbp]
\centering
\footnotesize
\setlength{\tabcolsep}{4pt}
\begin{tabular}{p{3.0cm}p{5.1cm}p{4.4cm}}
\toprule
Bloque final & Pregunta que resuelve & Evidencia requerida \\
\midrule
MAIN y controles directos & ¿Aprende el perfil congelado en la partición aleatoria y coinciden evaluación y validación independiente? & MAIN nueva, validación directa, \emph{bootstrap}, duplicados y etiquetas barajadas \\
Partición completa por días & ¿Persiste el comportamiento al romper la mezcla aleatoria entre días? & QRDQN y RF con lunes--miércoles para entrenamiento y jueves--viernes para prueba \\
Escalera de tamaños & ¿Cómo cambia el resultado con la cantidad de datos y el presupuesto proporcional de pasos? & Seis puntos lógicos entre $100\,000$ filas y el conjunto completo \\
Sensibilidad a semilla & ¿Cuánto varía el modelo al mantener fija la partición? & Cinco semillas de modelo, 42--46, sobre $1\,000\,000$ de filas \\
Escenarios retenidos & ¿Qué ataques o capturas inducen una degradación específica? & Cuatro \emph{holdouts} dirigidos y sus RF equivalentes \\
Comparación de modelos & ¿Qué parte del comportamiento depende de QRDQN y cuál del protocolo tabular? & Comparaciones emparejadas QRDQN--RF \\
Fase 2 & ¿Puede el MAIN final procesar la captura de laboratorio y con qué cambio de dominio? & Inferencia fresca, métricas y diagnósticos de distribución \\
\bottomrule
\end{tabular}
\caption[Estructura de los resultados finales]{Orden de incorporación de la campaña final y contrato mínimo de evidencia para cada bloque. La tabla describe decisiones analíticas, no resultados ejecutados. Fuente: elaboración propia a partir de \path{experiments/final_experiment_campaign.json}.}
\label{tab:estructura-resultados-finales}
\end{table}

\evidencenote{Los siete bloques de la Tabla~\ref{tab:estructura-resultados-finales} permanecen sin valores finales. La actualización posterior deberá producirse desde el agregado validado de campaña, nunca mediante transcripción estimada o reutilización silenciosa de cifras históricas.}

\section{Evidencia histórica previa a la campaña final}
\label{sec:evidencia-historica}

Los siguientes resultados proceden de artefactos comprometidos antes de congelar la campaña actual. Siguen siendo útiles para comprobar el funcionamiento del pipeline, identificar amenazas y justificar el diseño de la campaña; no deben leerse como una ejecución anticipada de esta.

\subsection{MAIN histórica: capacidad de aprendizaje intradistribución}
\label{sec:resultados-main}

La ejecución histórica \path{MAIN_qrdqn_cicids2017_canonical_full_random_20260609_193655} entrenó QRDQN durante $3\,000\,000$ de pasos sobre una partición aleatoria estratificada con semilla~42. Sobre $566\,149$ filas de prueba registró una exactitud de $0.99381$, \emph{recall} de ataque de $0.99536$ y F1 de ataque de $0.98445$. La matriz de confusión fue $\mathrm{TN}=451\,631$, $\mathrm{FP}=2\,989$, $\mathrm{FN}=518$ y $\mathrm{TP}=111\,011$.

\begin{table}[htbp]
\centering
\begin{tabular}{lc}
\toprule
Métrica & MAIN histórica \\
\midrule
Exactitud & $0.99381$ \\
Precisión de ataque & $0.97378$ \\
Recall de ataque & $0.99536$ \\
F1 de ataque & $0.98445$ \\
Tasa de falsos positivos & $0.00658$ \\
Tasa de falsos negativos & $0.00465$ \\
\bottomrule
\end{tabular}
\caption[Métricas de la MAIN histórica]{Métricas de la ejecución MAIN histórica sobre su partición aleatoria fija. No son resultados de la campaña final. Fuente: \texttt{metrics.json} y \texttt{bootstrap\_ci\_seed42.json} históricos.}
\label{tab:main-metrics}
\end{table}

Las curvas de la Figura~\ref{fig:curvas-entrenamiento} documentan que la ejecución completó un proceso de aprendizaje coherente; no permiten inferir por sí mismas generalización ni estabilidad entre semillas.

\begin{figure}[htbp]
\centering
\includegraphics[width=0.92\textwidth]{f8_curvas_entrenamiento.pdf}
\caption[Curvas de entrenamiento de la MAIN histórica]{Curvas de entrenamiento de la ejecución MAIN histórica. Fuente: escalares TensorBoard exportados del artefacto previo a campaña.}
\label{fig:curvas-entrenamiento}
\end{figure}

\begin{figure}[htbp]
\centering
\includegraphics[width=0.64\textwidth]{f9_cm_main.pdf}
\caption[Matriz de confusión de la MAIN histórica]{Matriz de confusión de la MAIN histórica sobre la partición aleatoria fija. Fuente: \texttt{bootstrap\_ci\_seed42.json}; no es la matriz de la MAIN final.}
\label{fig:cm-main}
\end{figure}

\subsection{Precisión de muestreo y duplicados históricos}
\label{sec:resultados-bootstrap}

El \emph{bootstrap} histórico de $10\,000$ remuestreos estratificados sitúa la F1 de ataque en $[0.98394,0.98496]$, el \emph{recall} de ataque en $[0.99495,0.99575]$, la tasa de falsos positivos en $[0.00634,0.00681]$ y la tasa de falsos negativos en $[0.00425,0.00505]$. Estos intervalos cuantifican la precisión sobre un conjunto de prueba fijo y un único modelo; no estiman varianza de entrenamiento.

\begin{figure}[htbp]
\centering
\includegraphics[width=0.88\textwidth]{f10_ic_bootstrap.pdf}
\caption[Intervalos bootstrap históricos]{Intervalos al 95\% de la MAIN histórica sobre su test fijo. Fuente: \texttt{runs/validation/bootstrap\_ci\_seed42.json}. La campaña final regenerará este análisis.}
\label{fig:ic-bootstrap}
\end{figure}

El análisis histórico de duplicados encontró que el $22.30\%$ de las filas completas eran duplicados exactos y que el $24.86\%$ del test aleatorio coincidía con alguna fila de entrenamiento. Entre los ataques de prueba, la coincidencia ascendía al $40.12\%$. No demuestra por sí solo una fuga de etiquetas en la implementación, pero sí que la mezcla aleatoria ofrece una evaluación optimista sobre un benchmark con observaciones repetidas.

\begin{figure}[htbp]
\centering
\includegraphics[width=0.84\textwidth]{f12_duplicados.pdf}
\caption[Duplicados en la partición aleatoria histórica]{Duplicados exactos y coincidencias entre entrenamiento y prueba de la partición histórica con semilla~42. Fuente: \texttt{runs/validation/duplicate\_analysis\_seed42.json}.}
\label{fig:duplicados}
\end{figure}

\subsection{Controles A, B y proxy temporal histórico}
\label{sec:resultados-escalera}

La Comprobación~A reprodujo la predicción directa del modelo sobre $10\,000$ muestras y confirmó que las métricas no dependían del bucle del entorno. La Comprobación~B reentrenó brevemente con etiquetas barajadas y obtuvo $0.4773$ de exactitud frente a $0.5227$ de la clase mayoritaria, señal compatible con la ausencia de una vía trivial de fuga en ese control.

La Comprobación~C entrenó una red proxy \texttt{[512,256]} durante $30\,000$ pasos con lunes--miércoles y evaluó jueves--viernes. Registró $0.52954$ de \emph{recall} y $0.62578$ de F1 de ataque. Esta ejecución es una señal histórica de dificultad temporal, pero no sustituye el experimento final: usa otra arquitectura, otro presupuesto y un contrato de artefactos anterior.

\begin{figure}[htbp]
\centering
\includegraphics[width=0.64\textwidth]{f11_cm_check_c.pdf}
\caption[Matriz histórica del proxy por días]{Matriz de confusión de la Comprobación~C histórica. Fuente: \texttt{VAL\_checks\_C\_20260213\_004847/validation\_results.json}. Se reemplazará como evidencia temporal principal por el experimento completo de $3\,000\,000$ de pasos.}
\label{fig:cm-check-c}
\end{figure}

\subsection{Random Forest histórico como referencia interpretativa}
\label{sec:rf}

El barrido histórico de \emph{Random Forest} empleó la misma representación canónica y particiones equivalentes. En la partición aleatoria superó a QRDQN histórica en F1 de ataque ($0.99676$ frente a $0.98445$). En la partición por día, su \emph{recall} de ataque descendió a $0.08135$ y su F1 a $0.15005$, frente a $0.52954$ y $0.62578$ del proxy QRDQN. Esta comparación motiva una evaluación emparejada actual, pero no demuestra que QRDQN generalice mejor de forma global: el QRDQN temporal histórico no era MAIN y faltan escenarios repetidos.

\begin{figure}[htbp]
\centering
\includegraphics[width=0.92\textwidth]{f13_qrdqn_vs_rf.pdf}
\caption[Comparación histórica entre QRDQN y Random Forest]{Recall y F1 de ataque en las particiones histórica aleatoria y por día. La fila temporal de QRDQN corresponde a un proxy; no existe aquí una comparación final de campaña. Fuente: artefactos históricos MAIN, Check~C y RF.}
\label{fig:qrdqn-vs-rf}
\end{figure}

\begin{figure}[htbp]
\centering
\includegraphics[width=0.64\textwidth]{f14_cm_rf_dia.pdf}
\caption[Matriz histórica de Random Forest por día]{Matriz de confusión de Random Forest en la partición histórica por día. Fuente: \texttt{rf\_cicids2017\_canonical\_20260628\_024735\_\_day\_split/metrics.json}.}
\label{fig:cm-rf-dia}
\end{figure}

\subsection{Inferencia histórica de Fase 2}
\label{sec:fase2}

La ejecución histórica \path{P2v2_pred_20260610_161231_MAIN} aplicó la MAIN histórica a $2\,000\,000$ de flujos etiquetados del laboratorio del autor. Registró $0.991862$ de exactitud, $0.988452$ de \emph{recall} de ataque, $0.98380$ de F1 de ataque y una tasa de bloqueo de $0.252364$. La captura era cerrada, generada por el operador y no adversarial; por ello el resultado solo acredita que ese pipeline procesó esa captura concreta. La inferencia fresca de campaña volverá a ejecutar el análisis con la MAIN final y sus artefactos de preprocesamiento.

\begin{figure}[htbp]
\centering
\includegraphics[width=0.80\textwidth]{f15_fase2.pdf}
\caption[Resultados históricos de Fase 2]{Matriz y tasa de bloqueo de la inferencia histórica sobre tráfico de laboratorio. Fuente: \texttt{runs/phase2/P2v2\_pred\_20260610\_161231\_MAIN/metrics.json}. Su validez externa es limitada.}
\label{fig:fase2-resultados}
\end{figure}

\section{Inserción de la campaña final}
\label{sec:resultados-finales-pendientes}

La actualización final de este capítulo seguirá la Tabla~\ref{tab:estructura-resultados-finales}. La MAIN fresca se presentará junto con su validación directa, \emph{bootstrap}, duplicados y control de etiquetas; después se analizarán la partición completa por días, la curva de tamaño, la distribución entre semillas y los cuatro escenarios retenidos. Solo entonces se realizará la comparación emparejada con Random Forest y se cerrará con la inferencia fresca de Fase~2. Este orden separa capacidad intradistribución, robustez y transferencia, y evita repetir la misma métrica en varias secciones.

\evidencenote{No existen todavía valores para la MAIN fresca, el día completo, la escalera de tamaños, las semillas 42--46, los cuatro holdouts, los RF nuevos, la validación directa, el bootstrap, los duplicados, el control barajado ni la Fase~2 fresca. Las figuras cuantitativas finales se generarán únicamente cuando el agregador de campaña declare el conjunto completo.}

\section{Síntesis provisional}
\label{sec:resultados-sintesis}

La evidencia histórica permite afirmar que la implementación aprende la tarea definida y que el protocolo aleatorio, aun siendo útil como comprobación interna, no basta para sostener generalización. También justifica mantener Random Forest como línea base fuerte y tratar la transferencia al laboratorio como una comprobación offline de alcance limitado. No permite determinar todavía la estabilidad entre semillas, la eficiencia respecto al tamaño, la sensibilidad de los cuatro escenarios, el comportamiento del perfil completo por días ni el balance final entre QRDQN y RF. Esas son, exactamente, las preguntas reservadas a la campaña final.

\chapter{Discusión}

La discusión no repite el inventario del Capítulo~\ref{sec:resultados-condiciones}; relaciona cada bloque de evidencia con la afirmación científica que puede sostener. Mientras la campaña final permanezca pendiente, los razonamientos que dependen de ella se expresan como criterios de decisión y no como resultados anticipados.

\section{Rendimiento intradistribución frente a generalización}

Una partición aleatoria estratificada responde a una pregunta necesaria pero limitada: si entrenamiento y prueba se extraen de la misma mezcla de CICIDS2017, ¿puede el pipeline aprender la regla binaria definida? La MAIN histórica responde afirmativamente para una ejecución anterior. Sin embargo, esta pregunta no equivale a saber si el modelo reconocerá ataques de otro día, captura o distribución. Los patrones de tráfico, herramientas, prevalencias y familias de ataque varían entre los CSV del benchmark, y una mezcla aleatoria diluye esa separación.

La campaña final trata la partición aleatoria como control intradistribución y reserva la afirmación de generalización al contraste con el experimento completo por días y los escenarios retenidos. Si MAIN mantiene métricas altas pero estas caen al separar días o CSV, la interpretación correcta será que el modelo explota bien la distribución mezclada, no que ofrece una defensa robusta. Si el comportamiento persiste, la afirmación podrá reforzarse únicamente para esos desplazamientos observados; seguirá sin equivaler a producción independiente.

\section{Duplicados y continuidad entre particiones}

El análisis histórico encontró una proporción relevante de coincidencias exactas entre entrenamiento y prueba. Esto no prueba una fuga de etiqueta dentro del código ni invalida automáticamente toda evaluación aleatoria: demuestra que el conjunto de prueba contiene menos novedad de la que su tamaño sugiere. En consecuencia, los intervalos estrechos obtenidos por \emph{bootstrap} pueden coexistir con una validez externa débil. El primero mide precisión sobre un test fijo; la segunda depende de la separación que ese test proporciona.

La regeneración final del análisis de duplicados debe interpretarse junto con el día completo. Si la tasa de coincidencias vuelve a ser elevada y el rendimiento temporal disminuye, ambas señales ofrecerán una explicación coherente del optimismo aleatorio. Si el rendimiento temporal se conserva, los duplicados seguirán siendo una amenaza del benchmark, pero no bastarán para atribuirles el resultado del modelo.

\section{Lectura sensible al coste}

La recompensa distingue cuatro consecuencias: permitir tráfico benigno, bloquear un ataque, bloquear tráfico benigno y permitir un ataque. La penalización más severa del falso negativo expresa una preferencia defensiva, pero no elimina el coste del falso positivo. Por eso la exactitud agregada no basta: deben leerse conjuntamente \emph{recall} de ataque, tasa de falsos negativos, precisión de ataque, tasa de falsos positivos y matriz de confusión.

Esta lectura es especialmente importante fuera de la partición aleatoria. Un modelo puede conservar una precisión alta de ataque porque bloquea muy pocos flujos y, a la vez, omitir la mayoría de las intrusiones. También puede mejorar \emph{recall} a costa de interrumpir tráfico legítimo. La campaña no fijará un umbral operativo ni ejecutará bloqueos reales; comparará el compromiso producido por las políticas entrenadas bajo una función de recompensa explícita.

\section{QRDQN frente a Random Forest}

Random Forest no se incluye como adversario débil, sino como prueba de necesidad. Con $\gamma=0$ y filas independientes, QRDQN resuelve una tarea próxima a clasificación tabular sensible al coste. Un clasificador supervisado sólido bajo el mismo esquema, partición y métrica permite preguntar qué aporta realmente la formulación distribucional, en lugar de atribuir a RL cualquier rendimiento que provenga de las características o del benchmark.

La evidencia histórica muestra que RF puede dominar la partición aleatoria y fallar ante una separación temporal concreta. Ese contraste es indicativo, pero está desbalanceado porque el QRDQN temporal histórico era un proxy. La campaña final corrige esta asimetría con pares equivalentes en MAIN, un millón de filas, día completo y cuatro escenarios retenidos. Si RF y QRDQN se degradan de forma parecida, la limitación principal apuntará a datos y dominio. Si difieren de forma consistente, podrá discutirse una sensibilidad específica de la familia de modelo. Un resultado favorable aislado no justificará superioridad general.

\section{Conjunto de datos como entorno y $\gamma=0$}

El uso de un conjunto tabular como entorno aporta una interfaz explícita de acciones y recompensas, compatibilidad con instrumentación de RL y una base para estudiar políticas sensibles al coste. A la vez, con $\gamma=0$ la devolución no incorpora consecuencias futuras: cada decisión depende de la observación actual y de su recompensa inmediata. La distribución cuantílica modelada por QRDQN describe la devolución de esa decisión, no una secuencia causal de estados de red.

Esta formulación enseña una lección metodológica: llamar \enquote{entorno} a un iterador sobre filas no convierte automáticamente la detección en un problema secuencial. La comparación supervisada es, por tanto, parte constitutiva de la evaluación. Una extensión genuinamente secuencial exigiría definir estado temporal, consecuencias de las acciones, episodios con dinámica de red y riesgos de intervención; queda fuera de este trabajo.

\section{Eficiencia de datos y sensibilidad a la semilla}

La escalera de tamaños distinguirá dos explicaciones que una única MAIN no puede separar: que el comportamiento dependa de casi todo CICIDS2017 o que se alcance con menos datos. Los presupuestos de pasos crecen proporcionalmente hasta el perfil completo, por lo que la curva debe leerse como respuesta conjunta a volumen y presupuesto de optimización. No medirá una ley universal de escalado.

El estudio de semillas responde a otra incertidumbre. Mantiene $1\,000\,000$ de filas y la semilla de partición~42, y cambia solo la semilla del modelo entre 42 y 46. Así estima sensibilidad de inicialización y aprendizaje sobre una partición fija; no incluye variación del conjunto de prueba. Media, dispersión, rango y casos atípicos deberán reportarse juntos. Una diferencia pequeña entre cinco ejecuciones reforzará la estabilidad local de ese punto, pero no sustituirá más particiones ni el MAIN completo repetido.

\section{Sensibilidad por escenario}

Los cuatro \emph{holdouts} ---Web Attacks, Infiltration, PortScan y DDoS--- fueron elegidos como escenarios dirigidos, no como barrido exhaustivo de los ocho CSV. La comparación debe mostrar el soporte y prevalencia de cada test junto con las métricas: un F1 no es comparable si la composición o la clase positiva cambian sustancialmente. El patrón entre escenarios será más informativo que una media única, porque permitirá identificar qué distribución retenida provoca cada error.

Esta selección limita la inferencia. Incluso si los cuatro escenarios fueran favorables, no quedaría demostrada invariancia frente a cualquier CSV, ataque o día. Si alguno falla, tampoco se deduce que la arquitectura sea inútil; sí se identifica una frontera concreta que requiere análisis de características, cobertura y posible cambio de dominio.

\section{Fase 2 y cambio de dominio}

La Fase~2 comprueba que el modelo, el escalador, los percentiles y el contrato canónico pueden aplicarse a flujos externos al cargador de CICIDS2017. Sus diagnósticos ---presencia de características, valores no finitos, recorte y puntuaciones estandarizadas--- ayudan a distinguir un error del pipeline de un desplazamiento de entrada. Cuando existan etiquetas, la matriz de confusión añade una señal de comportamiento.

No obstante, la captura pertenece a un único laboratorio cerrado, fue generada por el operador y no representa adversarios ni operación independiente. Un resultado alto demostraría compatibilidad y comportamiento sobre esa captura; uno bajo señalaría una limitación de transferencia o del mapeo. Ninguno de los dos, por sí solo, probaría preparación para despliegue.

\section{Amenazas a la validez y evidencia más fuerte}

La validez interna depende del contrato de partición, escalado solo con entrenamiento, semillas explícitas, control de etiquetas barajadas, validación directa y artefactos sellados. La validez de constructo está limitada por el mapeo binario y por tratar decisiones independientes como bandido contextual. La validez externa está limitada por CICIDS2017, cuatro escenarios seleccionados y una captura de laboratorio. La validez de conclusión estará condicionada por cinco semillas en un único tamaño y por la ausencia de repetición completa del MAIN.

Evidencia más fuerte requeriría capturas etiquetadas e independientes de varios entornos, repetición de los principales contrastes, comparación con más familias tabulares bajo el mismo contrato y evaluación previa a cualquier acción real. Esta tesis no pretende cubrir ese último nivel. Su objetivo es dejar una cadena reproducible y honesta que permita saber qué afirmación respalda cada experimento y cuál sigue fuera de alcance.

\chapter{Limitaciones}

El capítulo del protocolo experimental ya expuso las limitaciones \emph{conceptuales} del enfoque: el entorno estático, la formulación de bandido contextual sensible al coste, el mapeo de etiquetas binarias, la fragilidad del benchmark y el alcance de un único algoritmo. Este capítulo las retoma a la luz de los resultados medidos. La finalidad no es suavizar las conclusiones, sino dejar claro qué parte de la evidencia está respaldada por artefactos y qué parte queda todavía como trabajo pendiente.

\section{Modelo y semilla únicos}
Los resultados principales corresponden a un único modelo QRDQN entrenado con una sola semilla. La ejecución MAIN está bien documentada y sus artefactos permiten reproducir la lectura de métricas, pero no permite estimar la variabilidad propia del entrenamiento. En aprendizaje por refuerzo, incluso con la misma arquitectura y los mismos hiperparámetros, cambios de semilla pueden alterar la exploración, el orden efectivo de muestras, la inicialización de la red y la trayectoria de optimización. Esta tesis no mide esa variabilidad porque solo existe una ejecución oficial completa comprometida.

Los intervalos de confianza bootstrap reportados ($\pm 0.0002$ a $\pm 0.001$ al 95\%) cuantifican la precisión de muestreo del conjunto de prueba fijo, no la varianza entre ejecuciones de entrenamiento. Caracterizan cuán estable es la métrica sobre remuestreos del mismo conjunto de prueba para \emph{este} modelo, pero no cuánto variaría el rendimiento al reentrenar con semillas distintas. Las cifras deben leerse, por tanto, como el comportamiento de un punto de configuración bien caracterizado, no como una media sobre la aleatoriedad del entrenamiento.

\begin{table}[htbp]
\centering
\tableplaceholder{Varianza entre semillas de QRDQN: media, desviación e intervalo por métrica principal}
\caption[Tabla de varianza entre semillas pendiente]{\placeholdercaption{Varianza entre semillas de QRDQN: media, desviación e intervalo por métrica principal.}}
\label{tab:placeholder-varianza-semillas}
\end{table}

\section{Fuga en la partición aleatoria}
El análisis de duplicados muestra que el 22.30\% de las filas del conjunto completo son duplicados exactos y que el 24.86\% de las filas de prueba de la partición aleatoria duplican una fila de entrenamiento (40.12\% entre los ataques de prueba). Esto no significa que el modelo no aprenda nada: MAIN mantiene una matriz de confusión coherente y supera ampliamente una comprobación trivial. Sí significa, en cambio, que la partición aleatoria contiene información repetida que facilita la tarea y que sus métricas deben interpretarse como una cota optimista dentro del benchmark.

Esta limitación se mitiga metodológicamente reportando la Comprobación~C, basada en partición por día, como señal de generalización más exigente. Aun así, no se elimina por completo. La Comprobación~C no sustituye al experimento principal porque usa una red proxy \texttt{[512,256]} entrenada durante $30\,000$ pasos, no los pesos MAIN. Por tanto, la tesis conserva ambos peldaños con significados distintos: la partición aleatoria demuestra capacidad de aprendizaje bajo el protocolo principal; la partición por día muestra la fragilidad de la generalización cuando se rompe la continuidad del benchmark.

\section{Comparación con la línea base}
La comparación con Random Forest se ejecutó bajo el mismo protocolo: representación canónica, escalado, ponderación de clases balanceada y particiones equivalentes. Esto hace que la comparación sea mucho más fuerte que una comparación contra cifras externas de la literatura. En la partición aleatoria fija, Random Forest supera a QRDQN en F1 de ataque ($0.99676$ frente a $0.98445$), y esta ventaja debe declararse sin matices artificiales. La contribución de QRDQN no es dominar el benchmark cómodo, sino ofrecer una lectura sensible al coste bajo el cambio de partición por día.

Persisten dos limitaciones. La primera es que la comparación abarca un único modelo supervisado; otros clasificadores tabulares competitivos, como los ensamblados con gradient boosting, se discuten en la revisión de la literatura pero no forman parte del protocolo ejecutado. La segunda es que el barrido leave-one-CSV-out de Random Forest (miércoles reservado, recall de ataque $0.00719$) no tiene todavía un artefacto QRDQN comprometido equivalente. En ese peldaño, la comparación es solo de la línea base consigo misma y no entre modelos. La comparación entre QRDQN y Random Forest es directa en la partición aleatoria fija y en la partición por día, pero no en el leave-one-CSV-out.

\begin{table}[htbp]
\centering
\tableplaceholder{Comparación ampliada con modelos tabulares supervisados bajo el mismo protocolo}
\caption[Comparación ampliada de líneas base pendiente]{\placeholdercaption{Comparación ampliada con modelos tabulares supervisados bajo el mismo protocolo.}}
\label{tab:placeholder-lineas-base-ampliadas}
\end{table}

\section{Validez externa de la Fase 2}
El alto rendimiento de la Fase~2 ($0.991862$ de exactitud) se obtiene sobre tráfico generado por el operador en un laboratorio doméstico cerrado y no adversarial, de un único entorno y con etiquetas de intención del generador. No es tráfico de producción, no contiene un adversario real y no procede de una fuente independiente. Su validez externa es, por tanto, limitada: un buen resultado en la Fase~2 es una comprobación controlada de que el contrato de características y el preprocesamiento robusto transfieren a una captura cercana a la distribución, no una prueba de detección en el mundo real \parencite{SommerPaxson2010ClosedWorld,Layeghy2023CrossDomainNIDS}.

La principal aportación de la Fase~2 es de ingeniería experimental: demuestra que el pipeline puede procesar una captura externa al CSV de entrenamiento, aplicar el escalador persistido, conservar el esquema canónico y producir diagnósticos reproducibles. Ese resultado es necesario para defender que el sistema no vive únicamente dentro del cargador de CICIDS2017. Pero no basta para afirmar robustez operativa. Para ello harían falta capturas etiquetadas de varios entornos, condiciones benignas más diversas, tráfico adversarial independiente y un análisis de deriva más amplio.

\section{Experimentos planificados pero no ejecutados}
Por honestidad metodológica, se enumeran explícitamente los análisis previstos en el diseño cuyos artefactos no se han comprometido. El protocolo de curva de aprendizaje (escala de entrenamiento a distintos presupuestos) está diseñado e implementado, pero pendiente de GPU: no existe todavía ningún artefacto de curva de aprendizaje con métricas oficiales. La validación leave-one-CSV-out está implementada en el código para QRDQN, pero no se ha comprometido un artefacto agregado de sus métricas. La evaluación con múltiples semillas para estimar la varianza entre ejecuciones tampoco se ha llevado a cabo.

Ninguna de estas tres capacidades se reporta con cifras en esta tesis. Esta ausencia no es un detalle administrativo: afecta a la fuerza de las conclusiones. Sin curva de aprendizaje, no se puede cuantificar cómo cambia QRDQN al variar el presupuesto de entrenamiento; sin leave-one-CSV-out agregado, no se puede medir la estabilidad por escenario retenido; y sin múltiples semillas, no se puede separar el comportamiento del perfil MAIN de la varianza del entrenamiento. Declararlas como pendientes evita que el texto convierta un diseño experimental en evidencia ejecutada.

\begin{figure}[htbp]
\centering
\figureplaceholder{Escalera futura de validación externa: varios entornos etiquetados, tráfico independiente y análisis previo a despliegue activo}
\caption[Validación externa futura pendiente]{\placeholdercaption{Escalera futura de validación externa: varios entornos etiquetados, tráfico independiente y análisis previo a despliegue activo.}}
\label{fig:placeholder-validacion-externa}
\end{figure}

% F9.1 / Ch10.4 future-work stash:
% Las limitaciones anteriores definen una agenda de trabajo futuro coherente:
% ejecutar y comprometer la curva de aprendizaje y el leave-one-CSV-out de
% QRDQN; reentrenar con múltiples semillas para cuantificar la varianza;
% ampliar la comparación a otros modelos supervisados; y validar sobre tráfico
% etiquetado independiente de múltiples entornos antes de considerar cualquier
% paso hacia un despliegue activo, que a su vez exigiría un análisis de seguridad
% específico. Estas extensiones refuerzan la evidencia sin alterar la
% arquitectura experimental ya establecida.

\chapter{Consideraciones éticas y legales}

Las herramientas de detección de intrusiones operan sobre tráfico de red que puede contener información sensible, y un sistema que decide entre permitir y bloquear conlleva consecuencias operativas directas. Aunque esta tesis se limita a un prototipo experimental de inferencia offline, las decisiones de diseño que la hacen reproducible y segura tienen también una dimensión ética y legal que conviene hacer explícita.

\section{Uso dual y alcance estrictamente defensivo}
La investigación en ciberseguridad tiene una naturaleza intrínsecamente dual: las mismas técnicas que permiten detectar tráfico malicioso pueden, en principio, informar su evasión. El trabajo aquí descrito es defensivo por construcción y, además, deliberadamente inerte respecto a cualquier red real. Las salidas \texttt{PERMIT} y \texttt{BLOCK} son clasificaciones experimentales sobre registros de flujo estáticos: el sistema no realiza bloqueo en vivo, descarte de paquetes ni manipulación de \texttt{iptables}, y no se integra con cortafuegos, controladores de redes definidas por software (SDN) ni sistemas de prevención de intrusiones (IPS), tal como se delimita en los no objetivos del Capítulo~2. Esta ausencia de actuación en línea elimina el riesgo de que el artefacto, por sí mismo, provoque un bloqueo indebido de tráfico legítimo en una red en producción. El sistema es un instrumento de medición experimental, no un mecanismo de defensa desplegable.

\section{Privacidad y tratamiento de los datos}
El sistema opera exclusivamente sobre características estadísticas a nivel de flujo y no sobre el contenido (\emph{payload}) de los paquetes, evitando la inspección profunda de paquetes y buena parte de sus implicaciones legales y de privacidad. El contrato canónico de características refuerza esta posición al excluir a nivel de esquema los identificadores potencialmente sensibles o reidentificadores —direcciones IP de origen y destino, marcas temporales absolutas, \texttt{Flow ID} y números de puerto—, que se descartan antes de cualquier entrenamiento. El benchmark CICIDS2017 es un conjunto de datos público publicado por el Canadian Institute for Cybersecurity \parencite{Sharafaldin2018CICIDS2017,CICIDS2017Web}. El tráfico empleado en la Fase~2 fue generado por el propio operador en un laboratorio doméstico cerrado y no adversarial: se trata de tráfico propio, no contiene datos personales de terceros y no fue capturado de redes ajenas, por lo que su uso no plantea problemas de consentimiento ni de exposición de información personal.

\section{Riesgo operativo y asimetría de errores}
Las consecuencias de los dos tipos de error son profundamente asimétricas. Un falso negativo —un ataque clasificado como benigno y permitido— puede habilitar exfiltración de datos o interrupción del servicio; un falso positivo —tráfico benigno bloqueado— interrumpe comunicaciones legítimas y consume tiempo de análisis. El propio diseño de la función de recompensa codifica esta asimetría penalizando los falsos negativos con mayor severidad. Esta tesis reporta estos riesgos de forma honesta en lugar de ocultarlos tras la exactitud agregada: sobre la partición aleatoria fija \emph{in-distribution}, el modelo oficial MAIN presenta una tasa de falsos positivos de $0.00658$ y una tasa de falsos negativos de $0.00465$; sin embargo, bajo la partición por días, sustancialmente más exigente (Check~C), el recall de ataque desciende hasta $0.52954$. Esta caída del rendimiento ante un cambio de condiciones es precisamente la razón por la que el sistema se presenta como un prototipo de investigación y no como una herramienta lista para el despliegue: un uso operativo responsable exigiría validación sobre tráfico de producción etiquetado, múltiples entornos de captura y un análisis de seguridad independiente antes de delegar en el modelo cualquier decisión con efectos reales \parencite{Apruzzese2022CrossEvaluationNIDS,SommerPaxson2010ClosedWorld}.

\section{Reproducibilidad como práctica responsable}
Por último, la disciplina de reproducibilidad descrita en la metodología es también una práctica ética: la investigación en seguridad con aprendizaje automático adolece con frecuencia de código ausente y evaluaciones irreproducibles, lo que dificulta verificar o refutar las afirmaciones publicadas \parencite{Olszewski2023ReproSecurityML,Arp2020DosDontsMLSecurity}. La persistencia de artefactos por ejecución, la partición de prueba fija con semilla~42 y su manifiesto de referencia permiten que cualquier tercero audite exactamente bajo qué condiciones se obtuvo cada cifra reportada en esta tesis. Hacer las limitaciones verificables, en lugar de presentarlas como una declaración de buenas intenciones, es la forma de honestidad metodológica que este trabajo persigue.


\chapter{Bibliografía}
\printbibliography[heading=none]
\end{document}
